% ============================================================
% 第九章 ζ 函数的解析延拓与函数方程
% 内容：η 初等延拓；Riemann 第一法 (Hankel)；第二法 (ϑ)；平凡零点与临界带
% 目的：全平面亚纯延拓 + 函数方程；为 ξ 与零点章定点
% 编译方式：XeLaTeX（作为 main.tex 的 \input 子文件）
% ============================================================

\chapter{\texorpdfstring{$\zeta$}{Zeta} 函数的解析延拓与函数方程}\label{ch:riemann_zeta_continuation}

\section{本章导读}

\paragraph{定义域的阶段性}
此前各章对 $\zeta$ 的讨论主要限于\textbf{右半平面}
$\operatorname{Re}(s)>1$：Dirichlet 级数、Euler 乘积，以及第~\ref{ch:riemann_zeta_intro}~章中
借助 $\Gamma$（$\operatorname{Re}s>0$）写出的积分表示
$\zeta(s)=\frac1{\Gamma(s)}\int_0^\infty\frac{x^{s-1}}{e^x-1}\,dx$。
在该半平面外级数发散，故不能把上述表示直接当作全平面定义。
要理解 $\zeta$ 在 $\operatorname{Re}s\le 1$ 的形态（极点 $s=1$、临界带零点、与 $\Gamma$ 的函数方程），
必须扩大定义域——本章给出解析延拓与函数方程的具体构造。
（$\Gamma$ 自身已在第~\ref{ch:gamma_function}~章由 $\operatorname{Re}s>0$ 延拓至全平面；定义与唯一性见第~\ref{ch:analytic_continuation_general}~章。）

在 $\operatorname{Re}(s)>1$ 上，
\[
  \zeta(s) \;=\; \sum_{n=1}^{\infty} \frac{1}{n^s} \;=\; \prod_{p} \left(1 - p^{-s}\right)^{-1}.
\]
Dirichlet 级数在该半平面外发散，但 $\zeta$ 可亚纯延拓到 $\mathbb{C}$（仅 $s=1$ 为极点），并满足函数方程。
黎曼 1859 年论文给出两种经典构造（细节见正文与第~\ref{ch:riemann_original_interpretation}~章）：
\begin{itemize}
    \item \textbf{第一法（Hankel 围道）}：将积分表示转为绕正实轴分支切割的围道积分，再由留数得到 $\zeta(s)$ 与 $\zeta(1-s)$ 的关系。
    \item \textbf{第二法（$\vartheta$ 函数）}：利用 Jacobi $\vartheta$ 的模反演（定义及由 Gauss 型 Fourier 变换与 Poisson 求和给出的证明见第~\ref{ch:integral_transforms}~章），分离极点并得到全平面收敛的积分表示，从而导出对称形式的函数方程。
\end{itemize}

\textbf{本章安排}：先用 $\eta(s)$ \textbf{分步}初等延拓——
$\Re(s)>1\to$ 越过 $\Re(s)=1$（与第5章未延拓对照）$\to$ $\Re(s)>0$
（临界带/线与非平凡零点的数值显现；$\eta$\textbf{不}给 FE，\textbf{不}证 RH）；
再**分述**第一法、第二法；两法后**再比较**；
最后由函数方程\textbf{证明}平凡零点并\textbf{界定}临界带/临界线。

\textbf{要点}：函数方程的推导同时给出 $\zeta$ 的解析延拓。所得函数方程是 $\xi$ 的构造、零点区域限制与显式公式的出发点。

\paragraph{围道选择}
正文第一法使用 Hankel 围道处理 $(-z)^{s-1}$ 的分支；与矩形围道在合适变形下等价（见附录~\ref{app:contour}）。
选择 Hankel 是为与第~\ref{ch:gamma_function}~章余元公式的处理一致，并便于写分支切割。

% ============================================================
\section{基于交错级数 \texorpdfstring{$\eta(s)$}{eta(s)} 的初等延拓
（分步扩大定义域）}\label{sec:elementary_extension}

在给出第一法、第二法之前，先用 $\eta$ 作\textbf{分步}初等延拓，使读者看见定义域扩大时 $\zeta$ 的变化。
\begin{enumerate}
  \item \textbf{起点} $\Re(s)>1$：第~\ref{ch:riemann_zeta_intro}~章 Dirichlet 级数（已知）。
  \item \textbf{第一步} 越过 $\Re(s)=1$：与第5章\textbf{未延拓}截断曲面对照，说明解析延拓的必要。
  \item \textbf{第二步} 扩至 $\Re(s)>0$：临界带、临界线与非平凡零点的\textbf{数值}显现。
  \item \textbf{止步}：越不过 $\Re(s)=0$；$\eta$\textbf{推不出}函数方程（全平面与 FE 仍靠两法）。
\end{enumerate}

\subsection{交错级数 \texorpdfstring{$\eta(s)$}{eta(s)} 与延拓公式}

\begin{equation}
\label{eq:eta_definition}
  \eta(s) \;=\; \sum_{n=1}^{\infty} \frac{(-1)^{n-1}}{n^s}
  \;=\; 1 - \frac{1}{2^s} + \frac{1}{3^s} - \frac{1}{4^s} + \cdots.
\end{equation}
实数 $s>0$ 时为交错级数，收敛；复数 $\Re(s)>0$ 时\textbf{条件收敛}（Dirichlet 审敛），
\textbf{非}绝对收敛（绝对收敛 $\Leftrightarrow$ $\Re(s)>1$）。
在每个 $\{\Re(s)\ge\delta>0\}$ 的紧子集上一致收敛，故 $\eta$ 在 $\Re(s)>0$ 上全纯。

奇偶项分离（$\Re(s)>1$）得 $\zeta-\eta=2^{1-s}\zeta$，整理：
\begin{equation}
\label{eq:zeta_eta_continuation}
  \zeta(s) \;=\; \frac{\eta(s)}{1 - 2^{1-s}}.
\end{equation}
分母 $1-2^{1-s}=0$ 当且仅当
$s_k=1+\frac{2k\pi i}{\ln 2}$（$k\in\mathbb{Z}$），全在直线 $\Re(s)=1$ 上：
\begin{itemize}
  \item $k=0$：$s=1$。因 $\eta(1)=\ln 2\neq 0$，故 $s=1$ 为一阶极点；
    $\operatorname{Res}(\zeta,1)=\lim_{s\to1}(s-1)\zeta(s)=\eta(1)/\ln 2=1$
    （分母对 $s$ 求导用洛必达）。
  \item $k\neq0$：$s_k$ 为可去奇点，需 $\eta(s_k)=0$（Titchmarsh；
    或由后文全局延拓仅 $s=1$ 为极点反推）。补值后
    $F(s)=\eta(s)/(1-2^{1-s})$ 在 $\Re(s)>0\setminus\{1\}$ 全纯，且与 $\Re(s)>1$ 上的 $\zeta$ 一致。
\end{itemize}
边界上（$t\neq0$）写
\begin{equation}
  \zeta(1+it) \;=\; \frac{\eta(1+it)}{1 - 2^{-it}}.
  \label{eq:zeta_on_Re1}
\end{equation}

\subsection{第一步：越过 \texorpdfstring{$\Re(s)=1$}{Re(s)=1}
—— 与第~\ref{ch:riemann_zeta_intro}~章未延拓情形对照}

第~\ref{ch:riemann_zeta_intro}~章中 $\zeta$ 仅由 Dirichlet 级数定义在 $\Re(s)>1$。
图~\ref{fig:Modular}、\ref{fig:real}、\ref{fig:Imaginary} 上，级数曲面在 $\Re(s)=1$
\textbf{截断终止}：收敛边界如墙，墙右有定义，墙及左侧无定义。
从右侧逼近 $\Re(s)=1$（$t\neq0$）时，有限部分和剧烈振荡，边界在数值与几何上不连续。

若停留在「级数在哪收敛，函数就只在哪有定义」，则无法在 $\Re(s)=1$、$t\neq0$ 上谈论
$\zeta(1+it)$（素数定理所需边界），更无法进入临界带与零点理论。
故\textbf{解析延拓必要}：把 $\zeta$ 从「半平面上的级数」提升为「更大区域上的解析对象」。
用 \eqref{eq:zeta_eta_continuation}--\eqref{eq:zeta_on_Re1} 完成第一步：
$t\neq0$ 时定义推到直线 $\Re(s)=1$，除极点 $s=1$ 外边界有定义。

将图~\ref{fig:zeta_mod_half}--\ref{fig:zeta_imag_half} 与第5章未延拓三图\textbf{对照}：
\begin{itemize}
  \item \textbf{模}（图~\ref{fig:zeta_mod_half} $\leftrightarrow$ 图~\ref{fig:Modular}）：
        截断边缘消失，除 $s=1$ 模发散（一阶极点）外边界可衔接。
  \item \textbf{实部}（图~\ref{fig:zeta_real_half} $\leftrightarrow$ 图~\ref{fig:real}）：
        关于 $t=0$ 偶对称；$s=1$ 邻域左右符号跃迁与极点一致。
  \item \textbf{虚部}（图~\ref{fig:zeta_imag_half} $\leftrightarrow$ 图~\ref{fig:Imaginary}）：
        关于 $t=0$ 奇对称；$t=0$ 上虚部为零。
\end{itemize}
对照要点：\textbf{同一解析对象}越过级数收敛边界后仍存在；截断是表示法的局限，不是函数在边界外「消失」。

\begin{figure}[htbp]
\centering
\includegraphics[width=0.92\textwidth]{全书图形/8-Z模曲面-0.pdf}
\caption{第一步：$|\zeta|$ 推至 $\Re(s)=1$。对照图~\ref{fig:Modular}（级数截断）：
截断边缘消失，除 $s=1$ 极点外边界可衔接。红线示意推进位置。}
\label{fig:zeta_mod_half}
\end{figure}
\begin{figure}[htbp]
\centering
\includegraphics[width=0.92\textwidth]{全书图形/8-Z实曲面.pdf}
\caption{第一步：$\Re\zeta$。对照图~\ref{fig:real}；关于 $t=0$ 偶对称，$s=1$ 邻域左右发散方向相反。}
\label{fig:zeta_real_half}
\end{figure}
\begin{figure}[htbp]
\centering
\includegraphics[width=0.92\textwidth]{全书图形/8-Z虚曲面.pdf}
\caption{第一步：$\Im\zeta$。对照图~\ref{fig:Imaginary}；关于 $t=0$ 奇对称。}
\label{fig:zeta_imag_half}
\end{figure}

\subsection{第二步：扩至 \texorpdfstring{$\Re(s)>0$}{Re(s)>0}
—— 临界带、临界线与非平凡零点的数值显现}

在第一步已越过 $\Re(s)=1$ 的基础上，\eqref{eq:zeta_eta_continuation} 在整个开半平面
$\Re(s)>0$ 上有效，定义域向左逼近虚轴（\textbf{不含} $\Re(s)\le 0$）。
此时带状区域
\[
  0 < \Re(s) < 1
\]
已落在定义域内——即后文由函数方程\textbf{严格界定}的\textbf{临界带}的内部
（边界 $\Re(s)=0,1$ 上的无零点是更深结果，见第~\ref{ch:zero_distribution}~章；
此处先按开带作数值观察）。

平分该带的中轴线
\[
  \Re(s) = \tfrac12
\]
即图上的\textbf{临界线}（RH 断言非平凡零点均在此线上；本节\textbf{不}证明 RH）。

图~\ref{fig:zeta_modulus_cont_strip} 的模曲面上，可同时辨认：
\begin{enumerate}
  \item \textbf{一阶极点}：$s=1$ 处模发散，仍是右半平面内唯一极点。
  \item \textbf{临界带}：$0<\Re(s)<1$ 的竖直条带进入视野（相对第一步「只到 $\Re(s)=1$」，
        这是定义域\textbf{再向左一截}后才出现的结构）。
  \item \textbf{临界线}：带的中线 $\Re(s)=1/2$。
  \item \textbf{非平凡零点（数值）}：带内若干点模触底。它们是 $\eta$（或延拓后的 $\zeta$）的真实零点，
        称为\textbf{非平凡零点}，以区别于后文由函数方程得到的平凡零点 $s=-2,-4,\ldots$
        （平凡零点在左半平面，本阶段 $\eta$ \textbf{看不见}）。
        图中低零点落在 $\Re(s)=1/2$ 上，\textbf{与 RH 一致，但不代替证明}；
        也\textbf{不是}由 $\eta$ 推出的函数方程推论——$\eta$ 路径根本没有 FE。
\end{enumerate}

图~\ref{fig:zeta_modulus_cont_strip} 展示极点与带内触底并存；图为示意，\textbf{不代替}证明。
相对第一步：第一步解决「级数墙能否越过」；第二步解决「越过后再向左，临界带与非平凡零点如何\textbf{先在图上显现}」。
二者都是初等延拓的几何收获；\textbf{临界带/线的理论定义}与平凡零点的\textbf{证明}仍在两法 + 函数方程之后（第~\ref{sec:fe_geometry_trivial_zeros}~节）。

\begin{figure}[htbp]
\centering
\includegraphics[width=0.85\textwidth]{全书图形/8-Z模曲面-1.pdf}
\caption{第二步：$\Re(s)>0$ 上 $|\zeta|$。
$s=1$ 极点发散；临界带 $0<\Re(s)<1$ 内触底与已知低非平凡零点位置一致（数值观察）；
中轴线 $\Re(s)=1/2$ 为临界线（图示，非 RH 证明）。}
\label{fig:zeta_modulus_cont_strip}
\end{figure}

\subsection{止步：无函数方程；通向两法}

$\eta$ 在 $\Re(s)\le0$ 发散，\eqref{eq:zeta_eta_continuation} 失效；
关系中无 $s\leftrightarrow 1-s$，故\textbf{推不出函数方程}，也无左半平面与平凡零点。
第一步说明：\textbf{必须}延拓才能越过级数墙；
第二步说明：延拓到 $\Re(s)>0$ 后，临界带与非平凡零点可先被\textbf{看见}；
下一步（第一法、第二法）则在全平面完成延拓并建立函数方程，从而\textbf{证明}平凡零点并\textbf{界定}临界带——本章主线由此展开。

% ============================================================
\section{从初等延拓到全平面：黎曼两法概要}\label{sec:riemann_two_methods_blueprint}

前一节用交错级数 $\eta(s)$ 将 $\zeta(s)$ 延拓到 $\Re(s)>0$。
该方法在 $\Re(s)\le 0$ 失效，且\textbf{$\eta$ 推不出函数方程}。
为得到全平面亚纯延拓与函数方程，Riemann（1859）给出两种完整构造。
二者各走\textbf{独立的推导路径}，调用\textbf{不同的数学工具}，\textbf{互不作为对方的前提或引理}
（详细对照见两法正文之后的第~\ref{sec:riemann_two_methods_comparison}~节）。
本节只作\textbf{导读}（各是什么、后文在哪一节）。

\begin{enumerate}
    \item \textbf{第一法：围道积分与分支切割}
    （第~\ref{sec:riemann_first_method}~节）
    对多值因子作分支切割，构造包围正实轴的 Hankel 围道，将级数相关的实积分化为围道积分；再由留数得到 $\zeta(s)$ 与 $\zeta(1-s)$ 的关系。

    \item \textbf{第二法：$\vartheta$ 函数与模反演}
    （第~\ref{sec:theta_function_method}~节）
    利用 Jacobi $\vartheta$ 的模反演
    $\vartheta(1/x)=\sqrt{x}\,\vartheta(x)$
    （证明见第~\ref{ch:integral_transforms}~章），经 Mellin 型积分分离极点，得到全平面收敛表示，并自然通向 $\xi(s)$ 的对称。
\end{enumerate}

下面两节依次给出推导；读完后再看第~\ref{sec:riemann_two_methods_comparison}~节的对照。

% ============================================================
\section{黎曼第一法：围道积分与留数定理}\label{sec:riemann_first_method}

\subsection{从实积分开始：代数与连续的桥梁}

考虑 $\Gamma$ 函数的积分定义：对于 $\Re(s)>0$，
\[
  \Gamma(s) \;=\;  \int_0^{\infty} t^{s-1} e^{-t} \,\mathrm{d}t.
\]
由第~\ref{ch:gamma_function}~章缩放式~\eqref{eq:gamma-scaled}（取正整数 $n$），
\[
  \frac{1}{n^s}
  \;=\;
  \frac{1}{\Gamma(s)} \int_0^{\infty} x^{s-1} e^{-nx} \,\mathrm{d}x
  \qquad(\Re s>0).
\]
在 $\Re(s)>1$ 绝对收敛下，对 $n=1,2,\ldots$ 求和并交换求和与积分，得
\begin{equation}
\label{eq:zeta_integral_1}
  \zeta(s)
  \;=\;
  \frac{1}{\Gamma(s)} \int_0^{\infty} \frac{x^{s-1}}{e^x - 1} \,\mathrm{d}x.
\end{equation}
此即 $\zeta(s)$ 在 $\Re(s)>1$ 上的实积分表示（与第~\ref{ch:riemann_zeta_intro}~章~\eqref{eq:zeta-integral} 一致）。
当 $x\to 0$ 时被积函数 $\sim x^{s-2}$，故 $\Re(s)\le 1$ 时该实积分在原点附近发散；
为突破此限制，将路径变形为复围道。

\subsection{Hankel 围道与解析延拓}

在复平面上引入 Hankel 围道 $C$\footnote{关于 Hankel 围道与矩形围道的等价性，参见附录~\ref{app:contour}。}：

该围道从正实轴的 $+\infty$ 处出发，沿实轴上方（即 $z = x + i\epsilon$）向左平行移动，

在原点附近以逆时针方向绕圆周（半径为 $\epsilon$）一周，然后沿实轴下方（即 $z = x - i\epsilon$）平行返回至 $+\infty$。

为使多值函数 $(-z)^{s-1}=\exp\big((s-1)\log(-z)\big)$ 单值化，取分支切割线为正实轴 $[0,\infty)$。

在路径的实轴上方上沿 $z = x + i0$ 处，其辐角为 $\arg(-z) = -\pi$，故 $-z = x e^{-i\pi}$；

在路径的实轴下方下沿 $z = x - i0$ 处，其辐角为 $\arg(-z) = \pi$，故 $-z = x e^{i\pi}$。

综上有：
\[
  (-z)^{s-1} \;=\;
  \begin{cases}
    x^{s-1} e^{-i\pi(s-1)}, & \text{沿路径上半部分} \\
    x^{s-1} e^{i\pi(s-1)}, & \text{沿路径下半部分返回}
  \end{cases}
\]
\begin{figure}[htbp]
\centering
\begin{tikzpicture}[
    > = stealth,
    decoration={markings, mark=at position 0.55 with {\arrow[scale=1.5]{>}}}
]
    % 绘制坐标轴
    \draw[->, thick, gray] (-2.5, 0) -- (6.5, 0) node[right, text=black] {$\Re(z)$};
    \draw[->, thick, gray] (0, -2.5) -- (0, 2.5) node[above, text=black] {$\Im(z)$};

    % 绘制分支切割线
    \draw[ultra thick, brown, opacity=0.5, decorate, decoration={zigzag, segment length=2mm, amplitude=0.5mm}] (0,0) -- (5.8,0);
    \node[brown, below] at (4.5, -1) {分支切割线};

    % 绘制汉克尔围道
    \def\r{1.0}
    \def\eps{0.35}
    \def\R{5.8}
    \pgfmathsetmacro{\ang}{asin(\eps/\r)}

    \draw[thick, blue, postaction={decorate}] (\R, \eps) -- ({\r*cos(\ang)}, \eps) node[midway, above] {$z = x + i\epsilon$};
    \draw[thick, blue, postaction={decorate}] ({\r*cos(\ang)}, \eps) arc (\ang:360-\ang:\r);
    \draw[thick, blue, postaction={decorate}] ({\r*cos(\ang)}, -\eps) -- (\R, -\eps) node[midway, below] {$z = x - i\epsilon$};

    \node[blue, left] at (-\r-0.2, 0.7) {围道 $C$};

    % 原点标记
    \filldraw (0,0) circle (1pt) node[below right] {$O$};
\end{tikzpicture}
\caption{避开正半实轴分支切割线的汉克尔围道 示意图}
\label{fig:hankel}
\end{figure}

考虑如下复围道积分：
\begin{equation}
\label{eq:I_s}
  I(s) \;=\;  \int_C \frac{(-z)^{s-1}}{e^z - 1} \,\mathrm{d}z.
\end{equation}

\paragraph{围道积分的定义与解析延拓的两步法}
需要特别强调：积分 $I(s)$ 是\textbf{直接通过固定的 Hankel 围道 $C$ 定义的}，该围道不依赖于参数 $s$。

对于任意复数 $s$，被积函数在围道 $C$ 上（避开原点小圆弧）绝对收敛，因此 $I(s)$ 是复平面 $\mathbb{C}$ 上的\textbf{整函数}。

接下来的推导将分两步进行：
\begin{enumerate}
  \item \textbf{第一步（限制区域）}：
  在 $\Re(s)>1$ 的条件下，证明小圆弧积分趋于 $0$，从而将 $I(s)$ 简化为实轴上下沿路径的积分差，进而建立 $I(s)$ 与 $\zeta(s)$ 的关系：

  \[
    I(s) = -2i\sin(\pi s)\Gamma(s)\zeta(s),\qquad \Re(s)>1.
  \]
  \item \textbf{第二步（解析延拓）}：
  上式左端 $I(s)$ 是整函数，右端 $-2i\sin(\pi s)\Gamma(s)\zeta(s)$ 在 $\Re(s)>1$ 时等于 $I(s)$，因此该式给出了 $\zeta(s)$ 在 $\Re(s)>1$ 区域的表达式。

  由于右端因子在全平面解析（除极点外），此方程实现了 $\zeta(s)$ 到全复平面的解析延拓。
\end{enumerate}

此即“先在收敛域建立恒等式，再由左端延拓右端”的标准手法。

在 $\Re(s)>1$ 的假定下，当小圆弧半径 $\epsilon \to 0$ 时，评估围道 $C$ 在原点附近的积分贡献。

在半径为 $\epsilon$ 的圆弧段上，设 $z = \epsilon e^{i\phi}$（$\phi$ 从 $2\pi$ 减小到 $0$）。

分子模估计满足 $|(-z)^{s-1}| \le \epsilon^{\sigma-1} e^{\pi |\tau|}$，而分母满足 $|e^z - 1| \sim \epsilon$。

因此被积函数在圆弧段上的阶数为 $O(\epsilon^{\sigma-2})$。

圆弧的积分长度为 $2\pi \epsilon$，故圆弧积分的贡献受限于 $O(\epsilon^{\sigma-1})$。

当 $\Re(s) = \sigma > 1$ 时，随着 $\epsilon \to 0$，圆弧段的积分值确实收敛于 $0$。

此时，积分 $I(s)$ 简化为分别沿正实轴上、下两侧水平路径积分的代数和。

具体地，上沿路径（从 $+\infty$ 到 $0$）上 $-z = x e^{-i\pi}$，下沿路径（从 $0$ 到 $+\infty$）上 $-z = x e^{i\pi}$。

代入可得：
\begin{align*}
  I(s) &=  \int_{+\infty}^{0} \frac{x^{s-1} e^{-i\pi(s-1)}}{e^x - 1} \,\mathrm{d}x \;+\;  \int_{0}^{+\infty} \frac{x^{s-1} e^{i\pi(s-1)}}{e^x - 1} \,\mathrm{d}x \\
  &= \left( e^{i\pi(s-1)} - e^{-i\pi(s-1)} \right)  \int_0^{\infty} \frac{x^{s-1}}{e^x - 1} \,\mathrm{d}x \\
  &= 2i \sin\big(\pi(s-1)\big) \cdot \Gamma(s)\zeta(s) \\
  &= -2i \sin(\pi s) \Gamma(s) \zeta(s).
\end{align*}

其中最后一步利用了诱导公式：
\[
  \sin\!\big(\pi(s-1)\big) = \sin(\pi s - \pi) = -\sin(\pi s).
\]
此处利用当 $\Re(s)>1$ 时 $\zeta(s)$ 的积分表示公式 \eqref{eq:zeta_integral_1}。

由此得到 $\zeta(s)$ 在 $\Re(s)>1$ 条件下的围道积分表示：
\begin{equation}
\label{eq:analytic_cont_first}
  \zeta(s) \;=\; -\frac{1}{2i\sin(\pi s)\Gamma(s)}  \int_C \frac{(-z)^{s-1}}{e^z - 1} \,\mathrm{d}z.
\end{equation}

函数 $f(z)=\frac{(-z)^{s-1}}{e^z - 1}$ 除原点沿正实轴的分支切割线外，在复平面内其余部分存在单极点序列 $z=2n\pi i$。

鉴于汉克尔围道 $C$ 仅环绕原点而未包含其余极点，积分 $I(s)$ 对所有的复数 $s$ 均绝对收敛，因此 $I(s)$ 是 $s$ 的整函数。

结合 $\Gamma$ 函数的 Euler 反射公式 $\Gamma(s)\Gamma(1-s) = \frac{\pi}{\sin(\pi s)}$，上式可重写为：
\begin{equation}
\label{eq:zeta_contour_gamma}
  \zeta(s) \;=\; -\frac{\Gamma(1-s)}{2\pi i}  \int_C \frac{(-z)^{s-1}}{e^z - 1} \,\mathrm{d}z.
\end{equation}

由此将 $\zeta(s)$ 亚纯延拓到全平面（奇异性见下）。

为明确全平面上的奇异性，分析右端极点相消机制。

外侧的 $\Gamma(1-s)$ 在 $1-s = 0, -1, -2, \dots$ 处（即正整数 $s = 1, 2, 3, \dots$ 处）拥有简单的单极点。

然而，当 $s = n \in \{2, 3, 4, \dots\}$ 时，被积多值函数的指数 $s-1 = n-1$ 变为正整数。

此时，复平面的分枝切割线消失，被积函数 $\frac{(-z)^{n-1}}{e^z-1}$ 退化为单值亚纯函数。

由于 $n-1 \ge 1$，分子 $(-z)^{n-1}$ 在 $z=0$ 处的零点阶数至少为一阶，而分母 $e^z-1$ 在 $z=0$ 处为一阶零点，因此两零点相互抵消，使得 $z=0$ 实际上是被积函数的可去奇点（解析点）。

此时被积函数在全平面解析（除虚轴上的极点外），且在正实轴无割线。由于 Hankel 围道 $C$ 在无穷远闭合后不包含虚轴上的极点，根据 Cauchy 积分定理，该闭曲线上的积分值恒为 $0$，即：
\[
  I(n) \;=\;  \int_C \frac{(-z)^{n-1}}{e^z - 1} \,\mathrm{d}z \;=\; 0 \quad (n = 2, 3, 4, \dots).
\]
由此表明，整函数 $I(s)$ 在 $s = 2, 3, 4, \dots$ 处拥有至少一阶的零点，这恰好\textbf{完全抵消}了因子 $\Gamma(1-s)$ 在这些正整数点处的单极点。

而在 $s=1$ 处，被积函数为 $\frac{1}{e^z-1}$，$z=0$ 为其一阶极点且留数为 $1$。此时 Hankel 围道可收缩为仅环绕极点 $z=0$ 的闭曲线，其积分值为 $2\pi i \neq 0$，无法抵消 $\Gamma(1-s)$ 在 $s=1$ 处的极点。

因此，$\zeta(s)$ 在整个复平面内除 $s=1$ 处有一个留数为 $1$ 的简单单极点外，在其余所有点（包括所有正整数）均完全解析。

以上说明了极点相消机制。

\subsection{基于留数定理的函数方程推导}

进一步应用 Cauchy 留数定理。

考虑在复平面中引入由原围道 $C$ 与逆时针方向半径为 $R=(2N+1)\pi$ 的大圆周 $C_R$ 构成的环形闭合区域。

\subsubsection*{大圆周积分趋于零的条件与函数方程的有效区域}

设 $s=\sigma+i\tau$。在半径 $R=(2N+1)\pi$ 的大圆周 $C_R$ 上，有 $|-z| = R$。

分子部分满足模估计：
\[
  |(-z)^{s-1}| \;=\; |-z|^{\sigma-1} e^{-\tau \arg(-z)} \;\le\; R^{\sigma-1} e^{\pi |\tau|}.
\]
由于圆周的半径 $R = (2N+1)\pi$ 恰好穿过相邻两极点之间的中点，分母 $|e^z-1|$ 在整个圆周上存在正的下界，即 $|e^z-1| \ge c > 0$（其中常数 $c$ 与 $N$ 无关）。

因此，大圆周上的积分模长可以估计为：
\[
  \left| \oint_{C_R} \frac{(-z)^{s-1}}{e^z - 1} \,\mathrm{d}z \right| \;\le\; \frac{e^{\pi |\tau|}}{c} R^{\sigma-1} \cdot 2\pi R \;=\; \frac{2\pi e^{\pi |\tau|}}{c} R^{\sigma}.
\]
由上估计可得：
\begin{itemize}
  \item 当 $\sigma = \Re(s) < 0$ 时，$R^{\sigma} \to 0$（$R\to\infty$），大圆周积分严格趋于 $0$；
  \item 当 $\sigma = 0$ 时，$R^{\sigma} = 1$，积分可能不趋于 $0$（实际上有界但非零）；
  \item 当 $\sigma > 0$ 时，积分发散。
\end{itemize}

因此，利用留数定理将 $I(s)$ 表示为内部极点留数之和的步骤，仅在 $\Re(s)<0$ 时严格成立。

由此首先得到函数方程在 $\Re(s)<0$ 区域内的形式：
\[
\zeta(s) = 2^s(\pi)^{s-1}\sin\left(\frac{\pi s}{2}\right)\Gamma(1-s)\zeta(1-s),\qquad \Re(s)<0.
\]
由于右侧的 $\zeta(1-s)$ 在 $\Re(s)<0$ 时满足 $\Re(1-s)>1$，可由原始 Dirichlet 级数定义，因此该方程将 $\zeta(s)$ 从 $\Re(s)>1$ 的区域\textbf{直接延拓}到了 $\Re(s)<0$ 的区域。

对于带状区域 $0\le\Re(s)\le 1$ 内的值，可通过第二法积分表示（公式 \eqref{eq:analytic_continuation}）或直接由解析延拓的唯一性进行定义。

最终，函数方程通过解析延拓被证明对全平面 $s \in \mathbb{C} \setminus \{1\}$ 成立。

当 $\Re(s)<0$ 时，外边界大圆周上的积分趋于 $0$。

根据留数定理，围道 $C$ 上的积分值可以改写为内部分布在虚轴上的极点产生的留数之和。

被积函数 $f(z) = \frac{(-z)^{s-1}}{e^z - 1}$ 在整个复平面上仅在虚轴上的点 $z = 2n\pi i$ 与 $z = -2n\pi i$ （$n=1,2,\dots$）处存在一阶极点。

由于大圆周围道的走向，结果与极点留数和相差一个负号：
\[
   \int_{C} \frac{(-z)^{s-1}}{e^z - 1} \,\mathrm{d}z \;=\; -2\pi i \sum_{n=1}^{\infty} \left( \mathop{\mathrm{Res}}\limits_{z=2n\pi i} f(z) \;+\; \mathop{\mathrm{Res}}\limits_{z=-2n\pi i} f(z) \right).
\]
在计算极点处的留数时，须确定多值函数 $(-z)^{s-1} = \exp\big((s-1)\log(-z)\big)$ 在极点处的单值分支。

由于 Hankel 围道 $C$ 沿正实轴 $[0,\infty)$ 作分支切割，对任意 $z\in\mathbb{C}\setminus[0,\infty)$，将极坐标写为：
\[
  z \;=\; r e^{i\theta}, \quad 0 < \theta < 2\pi.
\]
相应地，复变量 $-z$ 的辐角 $\arg(-z)$ 必须在区间 $(-\pi, \pi)$ 内连续滑动，以保证其不跨越割线。

通过连续路径延拓确定 $-z$ 的解析分支：
\begin{enumerate}
  \item \textbf{上沿切口起步}：
  当 $z$ 位于实轴上侧趋近割线（$\theta\to 0^+$）时，定义 $\arg(-z) = -\pi$，此时 $-z = x e^{-i\pi}$；

  \item \textbf{下沿切口终局}：
  当 $z$ 逆时针绕过原点，最终移至实轴下侧趋近于割线时，即 $\theta \to 2\pi^-$,
  由于辐角随 $z$ 的逆时针运动连续增加了 $2\pi$，故此时 $\arg(-z) = -\pi + 2\pi = \pi$，对应 $-z = x e^{i\pi}$。
\end{enumerate}

根据这一辐角的连续演变规律，对于复平面切割区域内任意角度为 $\theta = \arg(z) \in (0, 2\pi)$ 的点，其 $-z$ 的辐角应通过沿围道内部的路径连续过渡确定，即满足：
\begin{equation}
\label{eq:argument_continuation}
  \arg(-z) \;=\; \theta - \pi \quad (\text{其中 } 0 < \theta < 2\pi).
\end{equation}

容易验证，当 $\theta \in (0, 2\pi)$ 时，$\arg(-z) = \theta - \pi$ 恰好落在单值化区间 $(-\pi, \pi)$ 内，从而保持了函数的全纯性与单值性。

现将该规律用于虚轴极点 $z = \pm 2n\pi i$：
\begin{itemize}
  \item \textbf{上半平面极点 $z = 2n\pi i$}：
  此极点位于正虚轴上，其辐角 $\theta = \frac{\pi}{2}$。
  由\eqref{eq:argument_continuation}，

  \[
    \arg(-z) \;=\; \frac{\pi}{2} - \pi \;=\; -\frac{\pi}{2}.
  \]
  因此 $-z$ 的取值为：

  \[
    -z \;=\; 2n\pi e^{-i\pi/2}.
  \]
  \item \textbf{下半平面极点 $z = -2n\pi i$}：
  该极点位于负虚轴上，其辐角 $\theta = \frac{3\pi}{2}$。
  同理，沿路径连续延拓后可得：

  \[
    \arg(-z) \;=\; \frac{3\pi}{2} - \pi \;=\; \frac{\pi}{2}.
  \]
  因此 $-z$ 的取值为：

  \[
    -z \;=\; 2n\pi e^{i\pi/2}.
  \]
\end{itemize}

由上述从围道边界向极点的连续延拓，辐角取值得以唯一确定，从而可计算各极点的留数。

在极点处求解分母导数（$e^z$），其值均为 $e^{\pm 2n\pi i} = 1$。

由此，上半平面极点处的留数为：
\[
  R^+ \;=\; \mathop{\mathrm{Res}}\limits_{z=2n\pi i} f(z) \;=\; \lim_{z \to 2n\pi i} (z - 2n\pi i) \frac{(-z)^{s-1}}{e^z - 1} \;=\; (2n\pi e^{-i\pi/2})^{s-1} \;=\; (2n\pi)^{s-1} e^{-i\pi(s-1)/2}.
\]
同理，下半平面极点处的留数为：
\[
  R^- \;=\; \mathop{\mathrm{Res}}\limits_{z=-2n\pi i} f(z) \;=\; (2n\pi e^{i\pi/2})^{s-1} \;=\; (2n\pi)^{s-1} e^{i\pi(s-1)/2}.
\]
将在此共轭极点对处的留数相加：
\begin{align*}
  R^+ + R^- &= (2n\pi)^{s-1} \left( e^{-i\pi(s-1)/2} + e^{i\pi(s-1)/2} \right) \\
  &= 2(2n\pi)^{s-1} \cos\left(\frac{\pi}{2}(s-1)\right) \\
  &= 2(2n\pi)^{s-1} \sin\left(\frac{\pi s}{2}\right).
\end{align*}

对所有正整数 $n$ 求和：
\[
  \sum_{n=1}^{\infty} 2(2n\pi)^{s-1} \sin\left(\frac{\pi s}{2}\right) \;=\; 2(2\pi)^{s-1} \sin\left(\frac{\pi s}{2}\right) \sum_{n=1}^{\infty} n^{s-1} \;=\; 2(2\pi)^{s-1} \sin\left(\frac{\pi s}{2}\right) \zeta(1-s).
\]
于是，原积分可表示为：
\[
  I(s) \;=\; -4\pi i (2\pi)^{s-1} \sin\left(\frac{\pi s}{2}\right) \zeta(1-s).
\]
因为 $I(s) = -2i \sin(\pi s)\Gamma(s)\zeta(s)$，比较等式两侧：
\[
  -2i \sin(\pi s)\Gamma(s)\zeta(s) \;=\; -4\pi i (2\pi)^{s-1} \sin\left(\frac{\pi s}{2}\right) \zeta(1-s).
\]
根据倍角公式 $\sin(\pi s) = 2 \sin\left(\frac{\pi s}{2}\right) \cos\left(\frac{\pi s}{2}\right)$ 简化，得出：
\begin{theorem}{黎曼函数方程}{riemann_functional_equation}
$\zeta(s)$ 在 $s\in\mathbb{C}\setminus\{1\}$ 上满足
\begin{equation}
\label{eq:asymmetric_fe}
  \zeta(1-s) \;=\; 2(2\pi)^{-s} \cos\left(\frac{\pi s}{2}\right) \Gamma(s) \zeta(s),
\end{equation}
以及与之等价的
\begin{equation}
\label{eq:asymmetric_fe_alt}
  \zeta(s) \;=\; 2^s\pi^{s-1} \sin\left(\frac{\pi s}{2}\right) \Gamma(1-s) \zeta(1-s).
\end{equation}
\end{theorem}

由 \eqref{eq:asymmetric_fe} 作替换 $s\mapsto 1-s$，得
\[
  \zeta(s) \;= \; 2(2\pi)^{-(1-s)} \cos\left(\frac{\pi(1-s)}{2}\right) \Gamma(1-s) \zeta(1-s).
\]
再用诱导公式
\[
  \cos\left(\frac{\pi(1-s)}{2}\right) \;=\; \cos\left(\frac{\pi}{2} - \frac{\pi s}{2}\right) \;=\; \sin\left(\frac{\pi s}{2}\right)
\]
整理指数项，即得 \eqref{eq:asymmetric_fe_alt}。故两种写法严格等价。

为阐明两端在 $s = n \in \{2, 3, 4, \dots\}$ 处的相消机制，除围道法中 $I(n)=0$ 外，亦可直接读函数方程两端的零极点：
\begin{itemize}
  \item 当 $s$ 为正偶数 $2k$ （$k \ge 1$）时，左侧 $\zeta(2k)$ 有限，右侧
        $\sin(\pi s/2)=\sin(k\pi)=0$ 的一阶零点与
        $\Gamma(1-s)=\Gamma(1-2k)$ 的一阶极点相消，乘积保持有限；
  \item 当 $s$ 为正奇数 $2k+1$ （$k \ge 1$）时，
        $\sin(\pi s/2)=(-1)^k\neq 0$，而 $\Gamma(1-s)=\Gamma(-2k)$ 有极点。
        由围道表示在正整数处 $I(n)=0$ 等已证机制
        （或本章稍后由函数方程推出的平凡零点 $\zeta(-2k)=0$），
        右侧相应因子相消，乘积保持有限。
        平凡零点 $\zeta(-2k)=0$ 是函数方程的\textbf{推论}而非先验假定；
        此处宜理解为两端亚纯延拓后的相容性检验，完整推导见下文“平凡零点”专节。
  \item 当 $s=1$ 时，$\sin(\pi/2)=1\neq 0$ 且 $\zeta(0)=-1/2\neq 0$，
        无零点消去 $\Gamma(1-s)$ 在 $s=1$ 的单极点，故 $s=1$ 为 $\zeta$ 的唯一一阶极点。
\end{itemize}
\begin{figure}[htbp]
\centering
\includegraphics[width=0.85\textwidth]{全书图形/8-函数方程.pdf}
\caption{解析延拓后 $\zeta(s)$ 与其反射项 $\zeta(1-s)$ 模曲面的双三维曲面叠加与反射对称几何形态（展示 $s=1$ 与 $s=0$ 处的对称发散峰顶，以及围绕中轴线 $\Re(s)=1/2$ 的交织起伏）}
\label{fig:zeta_global_surfaces}
\end{figure}

图 \ref{fig:zeta_global_surfaces} 给出 $\zeta(s)$ 模曲面与反射项 $\zeta(1-s)$ 模曲面的叠加示意，用以对照函数方程中的自变量反射 $s\mapsto 1-s$（图为示意，\textbf{不代替}证明）。图中可见：
\begin{enumerate}
    \item 模的叠加构成了一对对称的渐近峰顶结构。它们分别在各自的一阶极点 $s=1$（对应 $\zeta(s)$）与 $s=0$（对应 $\zeta(1-s)$）处展现出向无穷大发散的结构，并关于中轴线 $\Re(s) = 1/2$ 呈现出精确的反射对称，两者在此区域相互交织并向左侧振幅不断扩大地起伏。
    \item 由于在此阶段尚未对非平凡零点进行精细的理论分析与定位，图中未绘制零点标记，仅展示了由函数方程所决定的反射映射关系。
\end{enumerate}

% ============================================================
\section{黎曼第二法：\texorpdfstring{$\vartheta$}{theta} 函数与积分表示}\label{sec:theta_function_method}

本节给出与第一法\textbf{平行}的另一条构造：\textbf{不}引用上一节 Hankel 围道的中间结果，
而通过 $\vartheta$ 函数的模反演，在正实轴积分上建立关于 $s$ 与 $1-s$ 的对称结构。
Jacobi $\vartheta$ 的定义及其模反演已在第~\ref{ch:integral_transforms}~章给出，并由该章的 Gauss 型 Fourier 变换与 Poisson 求和完整证明；本节直接引用，重点转入由此导出的 $\zeta$ 积分表示。

\subsection{Jacobi \texorpdfstring{$\vartheta$}{theta} 函数与 \texorpdfstring{$\omega$}{omega} 反演}

第~\ref{ch:integral_transforms}~章已定义
\[
  \vartheta(x) \;=\; \sum_{n=-\infty}^{\infty} e^{-\pi n^2 x}, \qquad x>0,
\]
并借助 Gauss 型 Fourier 变换（第~\ref{subsec:gauss_ft}~节）与 Poisson 求和，证明了等价形式的模反演
\[
  \vartheta\!\left(\frac{1}{x}\right) \;=\; \sqrt{x}\, \vartheta(x)
  \quad\Longleftrightarrow\quad
  \vartheta(x) \;=\; \frac{1}{\sqrt{x}}\,\vartheta\!\left(\frac{1}{x}\right).
\]
下文不再重证，只由此转入第二法所需的 $\omega$ 形式。

在 $\Gamma$ 函数定义中考虑变量替换 $t = n^2 \pi x$：
\begin{equation}
\label{eq:gamma_sub_2}
  \pi^{-\frac{s}{2}} \, \Gamma\!\left(\frac{s}{2}\right) \frac{1}{n^s} \;=\;  \int_0^{\infty} x^{\frac{s}{2}-1} e^{-n^2 \pi x} \,\mathrm{d}x.
\end{equation}

在 $\Re(s)>1$ 条件下，两端关于 $n$ 从 $1$ 到 $\infty$ 求和：
\begin{equation}
\label{eq:zeta_integral_2}
  \pi^{-\frac{s}{2}} \, \Gamma\!\left(\frac{s}{2}\right) \zeta(s)
  \;=\;  \int_0^{\infty} x^{\frac{s}{2}-1} \left( \sum_{n=1}^{\infty} e^{-n^2 \pi x} \right) \,\mathrm{d}x.
\end{equation}

为了分析这一积分，定义函数 $\omega(x)$ ：
\[
  \omega(x) \;=\; \sum_{n=1}^{\infty} e^{-n^2 \pi x}.
\]
由第~\ref{ch:integral_transforms}~章的关系 $\vartheta(x)=1+2\sum_{n=1}^{\infty}e^{-\pi n^2 x}$，即
$\vartheta(x) = 1 + 2\omega(x)$。将模反演改写成 $\omega$ 的形式，得
\begin{equation}
\label{eq:omega_transform}
  \omega\!\left(\frac{1}{x}\right) \;=\; \frac{\sqrt{x}}{2} - \frac{1}{2} + \sqrt{x}\,\omega(x).
\end{equation}

\subsection{积分推演与极点分离}

回到式 \eqref{eq:zeta_integral_2} 的积分中。

为了解决 $x \to 0$ 端的收敛性，将积分区间划分为 $[0, 1]$ 和 $[1, \infty)$ 两部分：
\[
  \pi^{-\frac{s}{2}} \, \Gamma\!\left(\frac{s}{2}\right) \zeta(s) \;=\;  \int_0^1 x^{\frac{s}{2}-1} \omega(x) \,\mathrm{d}x \;+\;  \int_1^{\infty} x^{\frac{s}{2}-1} \omega(x) \,\mathrm{d}x.
\]
既然 $\omega(x)$ 在 $x \to \infty$ 具有指数衰减性质，区间 $[1, \infty)$ 上的积分对所有的复数 $s$ 均绝对收敛。

分析区间 $[0, 1]$ 上的第一部分积分（记为 $I_1$）：
\[
  I_1 \;=\;  \int_0^1 x^{\frac{s}{2}-1} \omega(x) \,\mathrm{d}x.
\]
作变量反演代换 $y = \frac{1}{x}, \mathrm{d}x = -y^{-2}\mathrm{d}y$：
\[
  I_1 \;=\;  \int_1^{\infty} y^{-\frac{s}{2}-1} \omega\!\left(\frac{1}{y}\right) \,\mathrm{d}y.
\]
将 $\omega(x)$ 的恒等式 \eqref{eq:omega_transform} 代入（为一致性，使用伪变量 $x$ 替代 $y$）：
\begin{align*}
  I_1
  &=  \int_1^{\infty} x^{-\frac{s}{2}-1} \left( \frac{x^{1/2}}{2} - \frac{1}{2} + x^{1/2}\omega(x) \right) \,\mathrm{d}x \\
  &= \frac{1}{2}  \int_1^{\infty} x^{-\frac{s+1}{2}} \,\mathrm{d}x \;-\; \frac{1}{2}  \int_1^{\infty} x^{-\frac{s}{2}-1} \,\mathrm{d}x \;+\;  \int_1^{\infty} x^{-\frac{s+1}{2}} \omega(x) \,\mathrm{d}x.
\end{align*}

前两项积分可以直接得出具有闭形式的结果：
\[
  \frac{1}{2} \left[ \frac{x^{\frac{1-s}{2}}}{\frac{1-s}{2}} \right]_1^{\infty} \;-\; \frac{1}{2} \left[ \frac{x^{-\frac{s}{2}}}{-\frac{s}{2}} \right]_1^{\infty} \;=\; \frac{-1}{1-s} \ - \ \frac{1}{s} \;=\; \frac{1}{s-1} \ - \ \frac{1}{s} \;=\; \frac{1}{s(s-1)}.
\]
这部分恰好构成了对复函数单极点特性的提取与分离。

\subsection{第二法下的解析延拓积分表示}

合并上述拆分的部分，可得：
\begin{align*}
  \pi^{-\frac{s}{2}} \Gamma\!\left(\frac{s}{2}\right) \zeta(s)
  &= \frac{1}{s(s-1)} \;+\;  \int_1^{\infty} x^{-\frac{s+1}{2}} \omega(x) \,\mathrm{d}x \;+\;  \int_1^{\infty} x^{\frac{s}{2}-1} \omega(x) \,\mathrm{d}x \\
  &= \frac{1}{s(s-1)} \;+\;  \int_1^{\infty} \left( x^{\frac{s}{2}-1} + x^{-\frac{s+1}{2}} \right) \omega(x) \,\mathrm{d}x.
\end{align*}
\begin{theorem}{$\zeta(s)$ 的积分解析延拓形式}{}
\begin{equation}
\label{eq:analytic_continuation}
  \pi^{-\frac{s}{2}} \, \Gamma\!\left(\frac{s}{2}\right) \zeta(s) \;=\; \frac{1}{s(s-1)} \;+\;  \int_1^{\infty} \left( x^{\frac{s}{2}-1} + x^{-\frac{s+1}{2}} \right) \omega(x) \,\mathrm{d}x, \quad s \in \mathbb{C} \setminus \{0,1\}.
\end{equation}

由于等式右端积分对全体复平面均绝对收敛，并且奇点部分只在 $s=0, 1$ 提供简单的代数极点，

该式提供了一个 $\zeta(s)$ 至复平面 $\mathbb{C}$ 上的亚纯函数解析延拓。
\end{theorem}

考察等式右端的形式。若作替换 $s \mapsto 1-s$：
\begin{itemize}
    \item 极点项 $\frac{1}{(1-s)(1-s-1)} = \frac{1}{s(s-1)}$ 表现为不变；
    \item 积分中的项 $\left( x^{\frac{s}{2}-1} + x^{-\frac{s+1}{2}} \right)$ 的两个组成部分互相映射交换位置，因此其和亦保持完全不变。
\end{itemize}

这表明，左端 $\pi^{-s/2} \Gamma(s/2) \zeta(s)$ 关于 $\Re(s)=1/2$ 具有反射对称性，为下一章定义满足 $\xi(s)=\xi(1-s)$ 的 $\xi(s)$ 作准备。

% ============================================================
\section{两法的比较与同一延拓}\label{sec:riemann_two_methods_comparison}
% ============================================================

第一法与第二法都已给出完整构造。下面分两点：\textbf{逻辑上是否独立、为何并列}；
\textbf{路径与产物有何不同、结果是否同一}。

\subsection{独立性与为何并列}

两法是两条\textbf{平行、自洽}的证明，而不是「主定理 + 另一写法」：

\begin{itemize}
  \item 各有\textbf{独立的推导路径}：第一法沿 Hankel 围道与留数推进；第二法沿 $\vartheta$ 模反演与全平面积分表示推进。
  \item 各调用\textbf{独立的数学工具包}：前者主要是分支、路径变形与留数；后者主要是 Gauss 型 Fourier 变换、Poisson 求和与模反演（第~\ref{ch:integral_transforms}~章），工具清单\textbf{不互相包含}。
  \item \textbf{互不作为对方的基础或引理}：第二法不引用第一法的中间结果，反之亦然；
        任取其一，即可单独得到全平面亚纯延拓与函数方程（中间公式形态可不同，见下表）。
\end{itemize}

正因为独立，书中\textbf{两法都写}才有信息量：不是同一步骤的复述，而是两条可独立检验的构造；
亦便于工具背景不同的读者择一细读，并与 Riemann 1859 两法并陈一致
（译介见第~\ref{ch:riemann_original_interpretation}~章）。

\subsection{路径区别与同一延拓}

表~\ref{tab:riemann_methods_comparison} 对照主路径、工具与后文衔接
（差别在\textbf{主演算画在何处}与直接产物形态；不作「复法 / 实法」对立）。

\begin{table}[htbp]
\centering
\begin{tabular}{p{2.2cm}|p{5.0cm}|p{5.0cm}}
\hline
\textbf{对比维度} & \textbf{第一法 (Hankel 围道)} & \textbf{第二法 ($\vartheta$ 模反演)} \\ \hline
\textbf{逻辑关系} & 自洽的完整构造；\textbf{不}依赖第二法
  & 自洽的完整构造；\textbf{不}依赖第一法 \\ \hline
\textbf{主路径} & 复平面围道、分支切割、留数
  & 正实轴上 $\vartheta$ / $\omega$、模反演，再 Mellin 型积分 \\ \hline
\textbf{主要工具} & 路径变形、分支、留数
  & Gauss 型 FT、Poisson 求和、模反演（第~\ref{ch:integral_transforms}~章） \\ \hline
\textbf{直接产物} & 围道积分表示；常得\textbf{非对称}型函数方程
  & 全平面收敛的积分表示 \eqref{eq:analytic_continuation}；$s\leftrightarrow 1-s$ \textbf{对称}更显 \\ \hline
\textbf{后文衔接} & 围道技术；与 1859 原稿围道叙述对照
  & 对称整函数 $\xi(s)$、Hadamard 乘积、临界线几何 \\ \hline
\end{tabular}
\caption{Riemann 解析延拓第一法与第二法对照（两法正文之后）}
\label{tab:riemann_methods_comparison}
\end{table}

路径与工具独立，并不产生两套 $\zeta$：二者均给出在 $\mathbb{C}\setminus\{1\}$ 上的亚纯延拓，
且在 $\operatorname{Re}s>1$ 与 Dirichlet 级数重合。由第~\ref{ch:analytic_continuation_general}~章
解析延拓的\textbf{唯一性}，这样的延拓至多一个，故两法定义的是\textbf{同一个}函数。

下一节由函数方程（任一法所得）导出临界带、临界线与平凡零点。

% ============================================================
\section{函数方程的推论：临界带、临界线与平凡零点}
\label{sec:fe_geometry_trivial_zeros}

两法已给出同一亚纯延拓（第~\ref{sec:riemann_two_methods_comparison}~节）。
本节由函数方程导出零点的全局区域限制，给出临界带与临界线的定义，并说明平凡零点的出现（图为示意，\textbf{不代替}证明）。
全平面延拓后的实部、虚部三维曲面及其零迹线，系统见第~\ref{ch:zero_distribution}~章。

\subsection{零点的全局区域限制与临界带、临界线的定义}
\label{sec:critical_strip_from_fe}

前文已建立非对称函数方程 \eqref{eq:asymmetric_fe_alt}。

据此可限制 $\zeta(s)$ 零点的全局分布区域。

$\zeta(s)$ 的零点分为两类：平凡零点，以及临界带内的非平凡零点。

非平凡零点的存在区域可如下确定：
\begin{enumerate}
    \item \textbf{右半平面的无零点区域}：
    当 $\Re(s) > 1$ 时，$\zeta(s)$ 可以表示为绝对收敛的 Euler 乘积：

    \[
      \zeta(s) \;=\; \prod_{p} \left(1 - p^{-s}\right)^{-1}.
    \]
    由于乘积中的每一项均不为零，且级数绝对收敛，因此在半平面 $\Re(s)>1$ 内 $\zeta(s)$ 恒不为零。

    \item \textbf{左半平面的零点限制}：
    当 $\Re(s) < 0$ 时，考虑非对称函数方程：

    \[
      \zeta(s) \;=\; 2^s \pi^{s-1} \sin\left(\frac{\pi s}{2}\right) \Gamma(1-s) \zeta(1-s).
    \]
    在右端项中：

    - 自变量反射项 $\zeta(1-s)$ 的实部满足 $\Re(1-s) > 1$，由前述结论可知 $\zeta(1-s) \neq 0$；

    - 因子 $\Gamma(1-s)$ 在全平面无零点（第\ref{ch:gamma_function}~章定理\ref{cor:gamma-no-zeros}），且在此处无极点；

    - 因此，等式右侧唯一的零点来源只能是正弦因子 $\sin\left(\frac{\pi s}{2}\right) = 0$。

    这直接导出了在负偶数处的\textbf{平凡零点}：

    \[
      s \;=\; -2, -4, -6, \dots
    \]
    也就是说，在左半平面 $\Re(s)<0$ 内，除了这些平凡零点外，不存在任何其他零点。
\end{enumerate}

综上所述，在仅用 $\operatorname{Re}s>1$ 无零点与函数方程时，所有非平凡零点必定被“挤压”在闭带
\[
  0 \;\leq\; \Re(s) \;\leq\; 1
\]
内。边界 $\Re(s)=1$（从而 $\Re(s)=0$）上的无零点是更深结果（Hadamard--de la Vallée Poussin，见第~\ref{ch:zero_distribution}~章），
得到开临界带 $0<\Re(s)<1$ 后，非平凡零点的精确定位区间才收紧。

据此给出：
\begin{definition}{临界带与临界线}{}
非平凡零点赖以存在的唯一带状区域 $0 \le \Re(s) \le 1$ 称为 $\zeta(s)$ 的\textbf{临界带}；

而平分该临界带的反射对称轴线 $\Re(s)=1/2$ 称为\textbf{临界线}。
\end{definition}

第~\ref{sec:elementary_extension}~节图~\ref{fig:zeta_modulus_cont_strip} 上所见的中轴线 $\Re(s)=1/2$，
在函数方程的对称性下正式获得作为「临界线」的定位（该图仍只是数值观察）。

\subsection{左半平面 \texorpdfstring{$\Re(s)<0$}{Re(s)<0} 上的平凡零点}
\label{sec:trivial_zeros_geometry}

前文用交错级数仅将定义域推至 $\Re(s)>0$；左半平面 $\Re(s)<0$ 在该阶段不可达。
完成全平面延拓后，可由函数方程讨论该半平面上的结构。

由函数方程，$\zeta(s)$ 在 $\Re(s)<0$ 上的行为由右半平面经 $\sin$、$\Gamma$ 与 $\zeta(1-s)$ 确定。
由非对称函数方程 \eqref{eq:asymmetric_fe_alt}，当 $\operatorname{Re}s\to-\infty$ 时，因子 $\Gamma(1-s)$ 与 $\sin(\pi s/2)$ 主导，左半平面曲面在虚部方向呈现振幅随 $|s|$ 增大而增长的振荡。

在\textbf{向左跨越虚轴}后，出现由函数方程直接给出的另一类零点——\textbf{平凡零点}（交错级数阶段不可见；上一小节已由区域限制推出其位置，此处补全论证与几何示意）：
\begin{itemize}
    \item 当 $s = -2, -4, -6, \dots$ 为负偶数（即 $s = -2n$, $n \in \mathbb{N}^+$）时，
    函数方程中的正弦因子 $\sin(\pi s/2) = \sin(-n\pi) = 0$ 恒为零。

    \item 此时 $\Gamma(1-s)=\Gamma(1+2n)$ 为有限正整数（\textbf{此处 $\Gamma(1-s)$ 无极点}，无需“极点--零点抵消”），
    且 $\zeta(1-s)=\zeta(1+2n)$ 亦为非零有限值，故正弦因子的零点直接给出：

    \[
      \zeta(-2n) \;=\; 0, \quad n = 1, 2, 3, \dots
    \]
\end{itemize}

图\ref{fig:trivial_zeros} 给出 $\zeta(s)$ 在左半实轴附近的局部曲面：在负偶数 $s=-2,-4,-6$ 处曲面与零平面相交，与由函数方程推出的平凡零点一致（图为示意，\textbf{不代替}证明；证明见上文）。
\begin{figure}[htbp]
\centering
\includegraphics[width=1.0\textwidth]{全书图形/8-Z模曲面-4.pdf}
\caption{$\zeta(s)$ 延拓至左半平面后的局部曲面图（实轴负偶数点 $s=-2,-4,-6$ 处曲面触零，与平凡零点一致；图为示意）}
\label{fig:trivial_zeros}
\end{figure}

全平面延拓后的实部、虚部三维曲面及其零迹线交点，以及与本章模曲面触底的对照，系统见第~\ref{ch:zero_distribution}~章第~\ref{sec:zeta_real_surface}、\ref{sec:zeta_imag_surface}~节。
$\xi(s)$ 的整函数化见第~\ref{ch:xi_function}~章；非平凡零点的几何与 $N(T)$ 亦见第~\ref{ch:zero_distribution}~章。

% ============================================================
\subsection{与第八章解析延拓定义与唯一性的联系}

第\ref{ch:analytic_continuation_general}章已说明：

\textbf{若一个函数满足某个函数方程，且该方程能将定义域从已知区域映射到新区域，则函数方程本身可作为解析延拓工具}

（参见第\ref{ch:analytic_continuation_general}章第\ref{subsec:functional_equation_method}节“函数方程”）。
\begin{theorem}{解析延拓的存在性与构造}{}
  由第一法或第二法推导出的黎曼函数方程，是 $\zeta(s)$ 从半平面 $\Re(s)>1$ 到全复平面 $s \in \mathbb{C} \setminus \{1\}$ 解析延拓的\textbf{显式构造}。

  特别地，第二法给出的积分表示（\eqref{eq:analytic_continuation}）和第一法给出的围道积分表示（\eqref{eq:analytic_cont_first}）均为 $\zeta(s)$ 在全平面上的有效定义。
\end{theorem}

这与第\ref{ch:analytic_continuation_general}章的函数方程延拓法一致，

同时也与 $\Gamma$ 函数利用递推公式 $\Gamma(s)=\Gamma(s+1)/s$ 向左延拓的做法同出一辙（参见第~\ref{ch:gamma_function}~章）。

至此，$\zeta(s)$ 的解析延拓不仅在理论上存在（由唯一性保证），而且通过函数方程得到了具体的、可计算的表达式。

% ============================================================
\section*{本章小结与后续展望}
\addcontentsline{toc}{section}{本章小结与后续展望}

本章给出 $\zeta$ 的两种全平面构造：先分述第一法（Hankel 围道与留数）、第二法（$\vartheta$ 模反演与积分表示）——
二者路径与工具\textbf{独立}，互不为引理；再在第~\ref{sec:riemann_two_methods_comparison}~节对照，
并说明由唯一性二者给出同一亚纯延拓（第~\ref{ch:analytic_continuation_general}~章）。

函数方程还给出关于临界线的对称，并推出左半平面上的平凡零点 $\zeta(-2k)=0$；
曲面图上可见 $s=1$ 极点与左半平面振荡（图为示意）。

$\zeta$ 在 $s=1$ 仍有极点，且函数方程中含 $\Gamma$ 与三角函数因子，
非平凡零点的乘积表示在 $\zeta$ 本身上不够干净。

下一章引入 $\xi(s)$，消去 $s=1$ 处的极点与函数方程中的附加因子，使对称写成
$\xi(s)=\xi(1-s)$，并为 Hadamard 乘积、零点计数与显式公式作准备。

\paragraph{知识谱系进度}
\label{sec:ch8_spectrum_progress}
本章完成了 $\zeta(s)$ 的解析延拓与函数方程。下图采用与第~\ref{ch:elementary_number_theoretic_functions}~章
图~\ref{fig:spectrum-stage2}、第~\ref{ch:spectrum_map}~章图~\ref{fig:function-relations}
\textbf{同一底图}（\texttt{黎曼猜想知识谱系关系图-A-2}）：
在图~\ref{fig:spectrum-stage2} 已着色的初等算术层之上，进一步着色 $\Gamma$、$\zeta$
及其延拓/函数方程相关连线（含 $\operatorname{Re}(s)>1$ 的 $\mu$、$\Lambda$、Euler 积分等分析桥梁）；
$\xi$、$N(T)$、Hadamard 乘积、RH，以及尚未在本章完成证明的\textbf{显式公式}连线
（$\zeta\to J$ 的 Perron 展开、$\zeta\to\psi$ 的 Perron 显式公式、$\psi\to\zeta$ 的 Mellin 负对数导）
仍保持灰白，留待第~\ref{ch:spectrum_map}~章全文着色。
读图时宜将\textbf{第二次着色新增彩色}元素与正文对照（表~\ref{tab:spectrum-stage8-guide}），
把图上的箭头读成已证公式的索引；初等层彩色元素仍见表~\ref{tab:spectrum-stage2-guide}。

% 图独占一页；题注文字排到下一页，避免与图例挤在同页
% （同一底图三次着色之第二次；\spectrumStage=8 为内部进度码，非章号）
\clearpage
\begin{figure}[p]
\centering
\def\spectrumStage{8}
\resizebox{\textwidth}{!}{%
% 黎曼猜想知识谱系关系图（A-2 现行底图 + 分阶段灰白）
% 用法：\def\spectrumStage{2|8|16}% 黎曼猜想知识谱系关系图（A-2 现行底图 + 分阶段灰白）
% 用法：\def\spectrumStage{2|8|16}% 黎曼猜想知识谱系关系图（A-2 现行底图 + 分阶段灰白）
% 用法：\def\spectrumStage{2|8|16}\input{图谱/谱系关系图-core.tex}
% 几何与 黎曼猜想知识谱系关系图-A-2.tex 同步；未涉及者仅灰白化（含图例）
%
% \spectrumStage 仅为内部进度码（历史命名），不是章号：
%   2  → 第一次着色：第2章 初等算术层（fig:spectrum-stage2）
%   8  → 第二次着色：第9章 延拓/函数方程（fig:spectrum-stage8）
%  16  → 第三次着色：第15章 全文总图（fig:function-relations）
% 正文表述用「第一次 / 第二次 / 第三次着色」或章标签，勿写「第8章图」「第16章图」。
\providecommand{\spectrumStage}{16}
\makeatletter
\colorlet{cDim}{gray!16}\colorlet{dDim}{gray!48}\colorlet{tDim}{gray!52}
\ifnum\spectrumStage=2
  \colorlet{cPrime}{blue!12}\colorlet{dPrime}{black}\colorlet{tPrime}{black}
  \colorlet{cMu}{teal!15}\colorlet{dMu}{teal!60!black}\colorlet{tMu}{black}
  \colorlet{cLi}{yellow!25}\colorlet{dLi}{orange!70!black}\colorlet{tLi}{black}
  \colorlet{cGammaU}{gray!16}\colorlet{cGammaL}{gray!16}\colorlet{dGamma}{gray!48}\colorlet{tGamma}{gray!52}
  \colorlet{cZetaU}{gray!16}\colorlet{cZetaL}{gray!16}\colorlet{tZeta}{gray!52}
  \colorlet{cXi}{gray!16}\colorlet{dXi}{gray!48}\colorlet{tXi}{gray!52}
  \colorlet{cNT}{gray!16}\colorlet{dNT}{gray!48}\colorlet{tNT}{gray!52}
  \colorlet{cProd}{gray!16}\colorlet{dProd}{gray!48}\colorlet{tProd}{gray!52}
  \colorlet{cRH}{gray!16}\colorlet{dRH}{gray!48}\colorlet{tRH}{gray!52}
  \colorlet{eArith}{blue!70!black}\colorlet{eElem}{violet!80!black}\colorlet{eDef}{black!80}
  \colorlet{ePNT}{brown!70!black}\colorlet{eAnal}{gray!42}\colorlet{eExpl}{gray!42}
  \colorlet{eExt}{gray!42}\colorlet{eZero}{gray!42}\colorlet{eLate}{gray!42}
  \colorlet{lgPrime}{blue!12}\colorlet{lgMu}{teal!15}\colorlet{lgLi}{yellow!25}
  \colorlet{lgXi}{gray!22}\colorlet{lgRH}{gray!22}\colorlet{lgGU}{gray!22}\colorlet{lgGL}{gray!22}
  \colorlet{lgEArith}{blue!70!black}\colorlet{lgEDef}{black!80}\colorlet{lgEElem}{violet!80!black}
  \colorlet{lgEZero}{gray!45}\colorlet{lgEAnal}{gray!45}\colorlet{lgEPNT}{brown!70!black}\colorlet{lgEExt}{gray!45}
  \colorlet{lgTprime}{black}\colorlet{lgTmu}{black}\colorlet{lgTli}{black}
  \colorlet{lgTxi}{gray!50}\colorlet{lgTrh}{gray!50}\colorlet{lgTgamma}{gray!50}
  \colorlet{lgTarith}{black}\colorlet{lgTdef}{black}\colorlet{lgTelem}{black}
  \colorlet{lgTzero}{gray!50}\colorlet{lgTanal}{gray!50}\colorlet{lgTpnt}{black}\colorlet{lgText}{gray!50}
\else
\ifnum\spectrumStage=8
  \colorlet{cPrime}{blue!12}\colorlet{dPrime}{black}\colorlet{tPrime}{black}
  \colorlet{cMu}{teal!15}\colorlet{dMu}{teal!60!black}\colorlet{tMu}{black}
  \colorlet{cLi}{yellow!25}\colorlet{dLi}{orange!70!black}\colorlet{tLi}{black}
  \colorlet{cGammaU}{blue!15}\colorlet{cGammaL}{red!15}\colorlet{dGamma}{black}\colorlet{tGamma}{black}
  \colorlet{cZetaU}{blue!15}\colorlet{cZetaL}{red!15}\colorlet{tZeta}{black}
  \colorlet{cXi}{gray!16}\colorlet{dXi}{gray!48}\colorlet{tXi}{gray!52}
  \colorlet{cNT}{gray!16}\colorlet{dNT}{gray!48}\colorlet{tNT}{gray!52}
  \colorlet{cProd}{gray!16}\colorlet{dProd}{gray!48}\colorlet{tProd}{gray!52}
  \colorlet{cRH}{gray!16}\colorlet{dRH}{gray!48}\colorlet{tRH}{gray!52}
  \colorlet{eArith}{blue!70!black}\colorlet{eElem}{violet!80!black}\colorlet{eDef}{black!80}
  \colorlet{ePNT}{brown!70!black}\colorlet{eAnal}{green!50!black}\colorlet{eExpl}{gray!42}
  \colorlet{eExt}{red!70!black}\colorlet{eZero}{gray!42}\colorlet{eLate}{gray!42}
  \colorlet{lgPrime}{blue!12}\colorlet{lgMu}{teal!15}\colorlet{lgLi}{yellow!25}
  \colorlet{lgXi}{gray!22}\colorlet{lgRH}{gray!22}\colorlet{lgGU}{blue!15}\colorlet{lgGL}{red!15}
  \colorlet{lgEArith}{blue!70!black}\colorlet{lgEDef}{black!80}\colorlet{lgEElem}{violet!80!black}
  \colorlet{lgEZero}{gray!45}\colorlet{lgEAnal}{green!50!black}\colorlet{lgEPNT}{brown!70!black}\colorlet{lgEExt}{red!70!black}
  \colorlet{lgTprime}{black}\colorlet{lgTmu}{black}\colorlet{lgTli}{black}
  \colorlet{lgTxi}{gray!50}\colorlet{lgTrh}{gray!50}\colorlet{lgTgamma}{black}
  \colorlet{lgTarith}{black}\colorlet{lgTdef}{black}\colorlet{lgTelem}{black}
  \colorlet{lgTzero}{gray!50}\colorlet{lgTanal}{black}\colorlet{lgTpnt}{black}\colorlet{lgText}{black}
\else
  \colorlet{cPrime}{blue!12}\colorlet{dPrime}{black}\colorlet{tPrime}{black}
  \colorlet{cMu}{teal!15}\colorlet{dMu}{teal!60!black}\colorlet{tMu}{black}
  \colorlet{cLi}{yellow!25}\colorlet{dLi}{orange!70!black}\colorlet{tLi}{black}
  \colorlet{cGammaU}{blue!15}\colorlet{cGammaL}{red!15}\colorlet{dGamma}{black}\colorlet{tGamma}{black}
  \colorlet{cZetaU}{blue!15}\colorlet{cZetaL}{red!15}\colorlet{tZeta}{black}
  \colorlet{cXi}{cyan!15}\colorlet{dXi}{cyan!60!black}\colorlet{tXi}{black}
  \colorlet{cNT}{cyan!15}\colorlet{dNT}{cyan!60!black}\colorlet{tNT}{black}
  \colorlet{cProd}{cyan!15}\colorlet{dProd}{cyan!60!black}\colorlet{tProd}{black}
  \colorlet{cRH}{orange!15}\colorlet{dRH}{orange!70!black}\colorlet{tRH}{black}
  \colorlet{eArith}{blue!70!black}\colorlet{eElem}{violet!80!black}\colorlet{eDef}{black!80}
  \colorlet{ePNT}{brown!70!black}\colorlet{eAnal}{green!50!black}\colorlet{eExpl}{green!50!black}
  \colorlet{eExt}{red!70!black}\colorlet{eZero}{orange!80!black}\colorlet{eLate}{black!80}
  \colorlet{lgPrime}{blue!12}\colorlet{lgMu}{teal!15}\colorlet{lgLi}{yellow!25}
  \colorlet{lgXi}{cyan!15}\colorlet{lgRH}{orange!15}\colorlet{lgGU}{blue!15}\colorlet{lgGL}{red!15}
  \colorlet{lgEArith}{blue!70!black}\colorlet{lgEDef}{black!80}\colorlet{lgEElem}{violet!80!black}
  \colorlet{lgEZero}{orange!80!black}\colorlet{lgEAnal}{green!50!black}\colorlet{lgEPNT}{brown!70!black}\colorlet{lgEExt}{red!70!black}
  \colorlet{lgTprime}{black}\colorlet{lgTmu}{black}\colorlet{lgTli}{black}
  \colorlet{lgTxi}{black}\colorlet{lgTrh}{black}\colorlet{lgTgamma}{black}
  \colorlet{lgTarith}{black}\colorlet{lgTdef}{black}\colorlet{lgTelem}{black}
  \colorlet{lgTzero}{black}\colorlet{lgTanal}{black}\colorlet{lgTpnt}{black}\colorlet{lgText}{black}
\fi\fi
\makeatother

\begin{tikzpicture}[scale=1.0, transform shape,
  box/.style    = {draw, rounded corners=5pt, minimum width=2.8cm,
                   minimum height=0.9cm, align=center, font=\small, thick},
  zbox/.style   = {draw, rounded corners=5pt, minimum width=2.8cm,
                   minimum height=0.9cm, align=center, font=\small,
                   thick, fill=gray!18},
  splitbox/.style = {draw, rounded corners=5pt, minimum width=7cm,
                     minimum height=2.2cm, align=center, font=\small, thick,
                     rectangle split, rectangle split parts=2,
                     rectangle split part fill={red!8, white}},
  arr/.style    = {-{Stealth[length=6pt,width=5pt]}, thick},
  barr/.style   = {{Stealth[length=6pt,width=5pt]}-{Stealth[length=6pt,width=5pt]}, thick},
  % 线型：虚线仅用于 N(T)→RH、ξ↔RH；其余为实线
  proven/.style = {solid},
  cond/.style   = {dashed},
  lab/.style    = {font=\tiny, align=center, inner sep=3pt},
  rhbox/.style  = {draw=dRH, rounded corners=5pt, minimum width=8cm,
                   minimum height=1.0cm, align=center, font=\small,
                   fill=cRH, text=tRH, thick}
]

% 布局与 A-2 同步：左移 1.0cm、上移 2.0cm
% 左右略放宽，避免 Γ/ζ 自环旁公式标注被裁切
\path[use as bounding box]
  (-10.2, -26.2) rectangle (10.0, 7.4);

\begin{scope}[shift={(-1.0,2.0)}]

\node[font=\bfseries\Large, align=center] at (0, 4.35)
  {黎曼猜想知识谱系关系图};

% 主图单独上移 1cm；标题与图例位置不动
\begin{scope}[shift={(-2.75,-1.2)}]
%% ==================== 节点分层 ====================

% 第一层：素数计数函数 (y=0)
\node[box, fill=cPrime, draw=dPrime]    (J)    at (-4.2,  0.0) {$J(x)$};
\node[box, fill=cPrime, draw=dPrime]    (pi)   at ( 0.0,  0.0) {$\pi(x)$};
\node[box, fill=cPrime, draw=dPrime]    (th)   at ( 5.0,  0.0) {$\theta(x)$};
\node[box, fill=cPrime, draw=dPrime]    (ps)   at ( 9.7,  0.0) {$\psi(x)$};

% 第二层：积性函数 (y=-3.0)
% 节点定义修改
\node[box, fill=cMu, draw=dMu] (mu)  at (0.3, -3.0) {$\mu(n)$};
\node[box, fill=cMu, draw=dMu] (Lam) at (5.0, -3.0) {$\Lambda(n)$};

%Weierstrass乘积
\node[box, fill=cProd, draw=dProd, text=tProd]  (prod)   at ( 9.5,  -17.0) {\tiny\text{Hadamard 乘积}};
% M(x) 节点已移除：初等等价信息直接标注在 RH 框内

% 第三层：Li 与 Gamma (y=-6.0 至 -7.0)
\node[box, fill=cLi, draw=dLi] (Li) at (-2.0, -6.0) {$\Li(x)$};

% Gamma(s) 上下分色：上部标收敛域，下部标延拓
\node[splitbox, rectangle split draw splits=false, minimum width=5cm,
      minimum height=4.7cm,
      rectangle split part fill={cGammaU, cGammaL}] (gamma) at (7.0, -6.0)
    {
     \centering \color{tGamma}$\operatorname{Re}(s)>0$
     \nodepart{two}
     \raisebox{6pt}{\color{tGamma}$\mathbb{C}\setminus\{0,-1,-2,\dots\}$}
    };
% Gamma(s) 标注（固定在节点中线左侧）
\node[font=\small, left=-30pt, text=tGamma] at (gamma.west) {$\Gamma(s)$};

% 【修订3】ζ(s) 节点：下半部分补全函数方程左侧 ζ(s)=
\node[splitbox, rectangle split draw splits=false, minimum height=4cm,
      rectangle split part fill={cZetaU, cZetaL}] (zeta) at (0.3, -10.5)
    {\centering\small
     \color{tZeta}$\displaystyle
       \sum_{n=1}^{\infty}\frac{1}{n^s}
       =\prod_{p}\dfrac{1}{1-p^{-s}}
       \quad(\operatorname{Re}(s)>1)$%
     \vphantom{$\displaystyle\sum_{n=1}^{\infty}$}%
     \nodepart{two}
     % 修订：补全 ζ(s)= 使函数方程完整
     \textcolor{tZeta}{$\displaystyle\zeta(s)=
       2^s\pi^{s-1}
       \sin\!\left(\frac{\pi s}{2}\right)\Gamma(1-s)\,\zeta(1-s)$%
     \vphantom{$\displaystyle\sum_{n=1}^{\infty}\mathbb{C}\setminus\{1\}$}}
    };
% ζ(s) 标注


% ξ(s) 节点：下移至与 RH 同一水平，右侧对称于 M(x)
\node[box, fill=cXi, draw=dXi, text=tXi, minimum height=1.2cm]
  (xi) at (4.85, -17.0)
  {$\xi(s)$\\ \tiny 非平凡零点$\rho$与$\zeta(s)$相同\\
  \tiny $\xi(s)=\xi(1-s)$(函数方程推论)
  };


% 【修订8】N(T) 节点：置于 ζ 与 RH 正中间，构成竖直逻辑链
\node[box, fill=cNT, draw=dNT, text=tNT, minimum height=1.3cm]
  (NT) at (2.3, -14.0)
  {$N(T)$\\ \tiny $\zeta(s)$ 零点计数函数\\[1pt]\tiny（幅角原理）};

% ==================== 黎曼猜想节点 ====================
% RH 下移至 y=-20.5，为 N(T) 留出充分空间
\node[rhbox, minimum width=4.6cm, minimum height=2.2cm] (rh) at (-1.9, -17.0) {
    {\tiny\color{tRH} 黎曼猜想 (Riemann Hypothesis, 1859)}\\[3pt]
    {\tiny\color{tRH}$\zeta(s)$ 的所有非平凡零点 $\rho$ 满足
    $\operatorname{Re}(\rho)=\dfrac{1}{2}$}\\[4pt]
{\tiny\color{tRH} 初等等价（梅滕斯函数）：%
  $\displaystyle M(x)=\sum_{n\leq x}\mu(n)$}\\
{\tiny\color{tRH}
  $\mathrm{RH} \iff M(x) = O(x^{1/2+\varepsilon}),\ \forall \varepsilon>0$}
};


% ==================== 连线 ====================

% ----- 第一层内部 -----
\draw[barr, eArith] (J.east) -- (pi.west)
  node[lab, midway, above=14pt, xshift=-2pt, eArith]
    {$\displaystyle\pi(x)=\sum_{n=1}^{\infty}\frac{\mu(n)}{n}J(x^{1/n})$}
  node[lab, midway, below=14pt, xshift=-2pt, eArith]
    {$\displaystyle J(x)=\sum_{n=1}^{\infty}\frac{1}{n}\pi(x^{1/n})$};

\draw[barr, eElem] (pi.east) -- (th.west)
  node[lab, midway, above=14pt, eElem]
    {$\displaystyle\theta(x)=\pi(x)\log x-\int_2^x\frac{\pi(t)}{t}\,\mathrm{d}t$}
  node[lab, midway, above={14pt}, yshift={-50pt}, eElem]
    {\tiny $\displaystyle\pi(x) = \dfrac{\theta(x)}{\log x}+ \int_{2}^{x} \frac{\theta(t)}{t\,(\log t)^2} \,\mathrm{d}t$};

\draw[barr, eElem] (th.east) -- (ps.west)
  node[lab, midway, above=11pt, xshift={5pt}, eElem]
    {$\displaystyle\psi(x)=\sum_{k=1}^{\infty}\theta(x^{1/k})$}
  node[lab, midway, above=-31pt, xshift={5pt}, yshift={-8pt}, eElem]
    {\tiny $\displaystyle\theta(x)=\sum_{k=1}^{\infty}\mu(k)\psi(x^{1/k})$};


% ----- 第二层内部 -----
\draw[barr, eArith] (mu.east) -- (Lam.west)
  node[lab, midway, yshift=20, xshift=0.3cm, eArith]
    {$\displaystyle\Lambda(n)=-\sum_{d\mid n}\mu(d)\log d$
     \quad{\rm (M\"{o}bius 反演)}}
  node[lab, midway, yshift={-12pt-0.5cm}, xshift=0, eArith]
    {$\displaystyle\log n=\sum_{d\mid n}\Lambda(d)$};

% μ(n) → M(x) 的定义连线已省去（见节点内文字），避免与已有连线交叉

% Lambda -> psi
\draw[arr, eDef] (Lam.east) to[out=0, in=-90, looseness=1.0]
  ($(ps.south west)+(30pt,0)$)
  node[lab, pos=0.5, yshift=-90, xshift=250, eDef]
    {$\displaystyle\psi(x)=\sum_{n\le x}\Lambda(n)$};

% ----- 第三层：Li 和 Gamma -----

% 【修订5】Li(x) → π(x)：改为渐近等价语义（非因果推导）
\draw[arr, ePNT]
  ($(Li.north)!.75!(Li.north east)$)
  to[out=90, in=-135, looseness=1.8]
  ($(pi.south west)+(35pt,0)$)
  node[lab, pos=0.4, yshift=-45, xshift=45, ePNT, align=left]
    {\textbf{素数定理（PNT）}：$\pi(x)\sim\Li(x)$\\
     (若RH成立，$\pi(x)=\Li(x)+O(\sqrt{x}\log x)$)};

% Li -> J（显式公式主项；属 eExpl：第2/9章灰白，第15章全文总图着色）
\draw[arr, eExpl]
  (Li.north) to[out=135, in=-90, looseness=1.5]
  ($(J.south east)+(-22pt,0)$)
  node[lab, pos=0.5, yshift=-115, xshift=-65, eExpl, align=center]
    {$\Li(x)$ 是 $J(x)$ 显式\\[2pt]公式的主项};

% Gamma（Re>0）→ ζ：Euler 积分表示（该积分公式本身要求 Re(s)>1）
\draw[arr, eAnal]
  ([yshift=6pt] gamma.west)
  to[out=180, in=90, looseness=0.8]
  ($(zeta.north)+(30pt,0)$)
  node[lab, pos=0.5, yshift=-235, xshift=90,eAnal, align=center]
    {$\displaystyle\zeta(s)\Gamma(s)
      =\int_0^{\infty}\frac{x^{s-1}}{e^x-1}\,\mathrm{d}x$\\
     $(\operatorname{Re}(s)>1)$};

% μ(n) → ζ(s)
\draw[arr, eAnal]
  (mu.south) to[out=-90, in=90, looseness=1.2]
  ($(zeta.north)+({-20pt+0.3cm},0)$)
  node[lab, pos=0.5, yshift={-120pt-0.5cm-0.4cm}, xshift={45pt-5pt}, eAnal]
    {$\displaystyle\frac{1}{\zeta(s)}=\sum_{n=1}^{\infty}\frac{\mu(n)}{n^s}$\\
     $(\operatorname{Re}(s)>1)$};

% ----- J(x) 与 ψ(x) 的 Stieltjes 积分互逆关系 -----
% 【修订4】加注 dψ(t)/ln t 在 t=1 附近需 t>1 的说明
\draw[barr, eArith, line width=0.8pt, rounded corners=8pt]
  ($(J.north)+(0,0.0cm)$)
  -- ++(0,2.0cm)
  -- ++(13.5cm,0)
  -- ++(0,-2.0cm)
  node[lab, pos=0.35, above=6pt, yshift={10pt+0.5cm}, xshift=-200, eArith, align=center]
    {$\displaystyle\psi(x)=\int_{2^-}^x\log t\,\mathrm{d}J(t)$\quad(Stieltjes 积分，$J\to\psi$)}
  node[lab, pos=0.65, below=6pt, yshift={32pt+0.2cm}, xshift=-205, eArith, align=center]
    {$\displaystyle J(x)=\int_{2^-}^x\frac{\mathrm{d}\psi(t)}{\log t}$
     \quad(Stieltjes 反演，$\psi\to J$；积分自 $t>1$)};

% ----- 第四层：ζ 和 ξ 与第二层的连线 -----

% ζ → ψ：显式公式（Perron 积分 + 留数，单向；x≠p^k 时公式精确成立）
\draw[arr, eExpl, rounded corners=30pt]
  ($(zeta.east)+(0, -10pt)$) -| ($(ps.south)+(30pt,0)$)
  node[lab, pos=0.85, right=10pt, yshift=-170, xshift={-170pt-0.2cm},eExpl, align=left]
{$\psi(x)$ 显式公式（Perron 积分，$\zeta\to\psi$）\\
     $\displaystyle \psi(x) = \frac{1}{2\pi i}\int_{c-i\infty}^{c+i\infty}\!\left(-\frac{\zeta'(s)}{\zeta(s)}\right)\!\frac{x^s}{s}\,\mathrm{d}s$\\
     $\displaystyle = x-\sum_{\rho}\frac{x^{\rho}}{\rho}-\log(2\pi)-\tfrac{1}{2}\log(1-x^{-2})$\\
     \tiny $x>1,\;x\neq p^k$；$\rho$ 按 $|\mathrm{Im}(\rho)|$ 升序配对，条件收敛};

% ψ → ζ：Mellin 变换（单向）
\draw[arr, eExpl, rounded corners=30pt]
  ($(ps.south)+(40pt,0)$) |- ($(zeta.east)+(0, -20pt)$)
  node[lab, pos=0.8, yshift=-27, xshift=30,eExpl, align=center]
    {Mellin 变换（$\psi(x) \longleftrightarrow -\dfrac{\zeta'(s)}{\zeta(s)}$，即负对数导数）\\
     $\displaystyle -\frac{\zeta'(s)}{\zeta(s)} = s\int_{1}^{\infty}\frac{\psi(x)}{x^{s+1}}\,\mathrm{d}x$};
       
       
       
       
       
       
       
% Lambda -> zeta（对数微分 Dirichlet 级数）
\draw[arr, eAnal]
  ($(Lam.south west)+(6pt,0)$) to[out=-90, in=90, looseness=0.4]
  ($(zeta.north west)!0.50!(zeta.north east)$)
  node[lab, pos=0.5, yshift={-120pt-0.5cm-0.4cm+1cm}, xshift={145pt-5pt}, eAnal]
    {$\displaystyle-\frac{\zeta'(s)}{\zeta(s)}
      =\sum_{n=1}^{\infty}\frac{\Lambda(n)}{n^s}$};

% J(x) 显式公式（ζ → J）
\draw[arr, eExpl, rounded corners=50pt]
  ($(zeta.south west)+(0,6pt)$) -- ++(-60pt,0) -- ($(J.south west)+(5pt,0)$)
  node[lab, pos=0.5, yshift=-203, xshift={75pt-0.2cm}, eExpl, align=center]
{\textbf{Perron 公式} $\Rightarrow$ $J(x)$ 显式展开\\
     $\displaystyle
     \begin{aligned}
         &J(x)=\frac{1}{2\pi i}\int_{c-i\infty}^{c+i\infty}
           \frac{\log\zeta(s)}{s}\,x^s\,\mathrm{d}s\\
         &=\Li(x)-\sum_{\rho}\Li(x^{\rho})
           -\log 2+\int_x^{\infty}\frac{\mathrm{d}t}{t(t^2-1)\log t}
     \end{aligned}$\\
     \tiny $x>1,\;x\neq p^k$；$\rho$ 按 $|\mathrm{Im}(\rho)|$ 升序配对，条件收敛};

% ζ → J：Mellin 变换（ln ζ）
\draw[arr, eAnal, rounded corners=5pt]
  (zeta.north west) to[out=110, in=-90, looseness=1.5] (J.south)
  node[lab, pos=0.4, yshift=-205, xshift=-55, eAnal, align=center]
    {$\displaystyle\log\zeta(s)=s\int_1^{\infty}J(x)\,x^{-s-1}\,\mathrm{d}x$\\
     $(\operatorname{Re}(s)>1)$};

% ζ → ξ（定义）：ξ 已下移，连线沿右侧空白区域向右下延伸
\draw[arr, eLate, rounded corners=15pt]
  ($(zeta.south east)!0.15!(zeta.north east)$)
  -- ++(1.0, 0)
  -- ++(0, -3.8)
  -- ($(xi.north)+(0pt,0)$)
  node[lab, pos=0.25, right=4pt,yshift=30, xshift=0, eLate, align=center]
    {$\xi(s)=\dfrac{1}{2}s(s-1)\pi^{-s/2}\Gamma(s/2)\zeta(s)$  
};






% xi -> prod (Weierstrass product connection)
\draw[arr, eLate]
  %(xi.south) -- (prod.south)
  (xi.east) --(prod.west)
  node[midway, right=4pt, yshift=28, xshift=-45, eLate]
    {\tiny$\displaystyle\xi(s)=e^{A+Bs}\prod_{\rho}\left(1-\frac{s}{\rho}\right)e^{s/\rho}$};

% prod -> ps (zero sum in explicit formula)
%\draw[arr, violet!70!black, rounded corners=15pt]
  %(prod.east) -- ++(4.4, 0) |- (ps.south)
  %node[lab, pos=0.6, right=2pt, violet!70!black, align=left]
   % {对数微分与 Perron 积分\\
    % 给出显式公式中\\
    % 零点和 $\sum_{\rho}\frac{x_{\rho}}{\rho}$ 项};
%
% ζ 自身左侧弧线（解析延拓）：与表 tab:spectrum-stage8-guide 一致
\draw[arr, eExt, rounded corners=5pt]
  ($(zeta.north west)!0.15!(zeta.south west)$)
  to[out=180, in=180, looseness=1.6]
  ($(zeta.north west)!0.85!(zeta.south west)$);
% 标签用独立节点锚定在 ζ 西侧弧线旁（避免 path 末 node 漂移）
% 单行；中心在 zeta.west 右 2.5cm（相对初版再累计右移）
\node[lab, eExt, align=center, anchor=center]
  at ([xshift=2.5cm]zeta.west)
  {$\zeta$ 亚纯延拓至 $\mathbb{C}\setminus\{1\}$（Hankel / $\vartheta$ 两法）};

% ζ → N(T)：竖直向下（弧线）
\draw[arr, eZero]
  ($(zeta.south east)!0.10!(zeta.north east)$)
  to[out=0, in=90] ++(0.8, -0.8)
  to[out=-90, in=0] (NT.east)
  node[lab, pos=0.6, right=6pt, yshift=-357, xshift=5, eZero, align=center]
    {$\displaystyle N(T)=\frac{T}{2\pi}\log\frac{T}{2\pi e}+O(\log T)$\\
     \tiny 零点计数公式（von Mangoldt）};

% 【已删去】ζ → RH 蓝色直连线（设计意图改为ζ→N(T)→RH线性链）
%\draw[arr, eArith, rounded corners=10pt]
%  (zeta.south)
%  -- ++(0, -3.35)
%  -- ++(-4.1, 0)
%  -- ($(rh.north)!0.80!(rh.north west)$)
%  node[lab, pos=0.35, left=6pt, yshift=20pt,xshift=120,eArith, align=center]
%    {
%     $\zeta(s)$ 的非平凡零点全部位于直线 $\operatorname{Re}(s)=\dfrac12$ 上};
% RH -- P(s) 画线 + 文字
%\draw[arr, cyan!60!black]
 % ($(rh.east)!0.60!(rh.south east)$) -- (prod.west)    % 起点 -- 终点（必须有！）
  %node[midway, right=6pt, font=\tiny, cyan!60!black, align=center]
  %  {$\zeta(s) = \frac{e^{s(\log(2\pi) - 1)}}{2(s-1)} \prod_{n=1}^{\infty} \left(1 + \frac{s}{2n}\right)e^{-s/2n} P(s)$};

% ξ(s) ↔ RH、N(T) → RH：虚线标注（与 RH 相关的谱系连线）

% N(T) → RH：完成 ζ→N(T)→RH 线性逻辑链（虚线）
\draw[arr, eZero, cond, rounded corners=8pt]
  (NT.south west)
  to[out=-150, in=60, looseness=0.9]
  ($(rh.north)+(20pt, 0)$)
node[lab, pos=0.55, left=4pt, yshift=-430, xshift=-25, eZero, align=center]
    {$\mathrm{RH} \iff N_0(T)=N(T),\ \forall\,T>0$\\
   \tiny（$N_0(T)$：临界线 $\operatorname{Re}(s)=\dfrac{1}{2}$ 上虚部 $\leq T$ 的零点数）};

\draw[barr, eZero, cond]
  (xi.west) -- ($(rh.east)!0.0!(rh.south east)$)
  node[lab, midway, above=20pt,yshift=3pt,xshift=5, eZero, align=center]
    {$\mathrm{RH}\iff\xi\!\left(\dfrac{1}{2}+it\right)$\\
     的所有零点 $t$ 均为实数};

% M(x)→RH 箭头已删去：等价信息直接写在 RH 框内

% Gamma 自身解析延拓弧线：与表 tab:spectrum-stage8-guide 一致
\draw[arr, eExt, rounded corners=1.5pt]
  ($(gamma.north east)+(0,-5pt)$)
  to[out=0, in=0, looseness=1.6]
  ($(gamma.south east)+(0,5pt)$);
% 标签用独立节点锚定在 Γ 东侧弧线旁（仅递推式；延拓域已在框下半）
% 相对初版 (gamma.east+14pt)：净向左 3.7cm，上移 2pt
\node[lab, eExt, align=center, anchor=west]
  at ($(gamma.east)+({14pt-3.7cm},2pt)$)
  {递推延拓 $\Gamma(s)=\Gamma(s+1)/s$};

% Gamma（延拓后）→ ζ：Hankel 围道
\draw[arr, eExt]
  ($(gamma.south west)+(40pt,0)$)
  to[out=180, in=0, looseness=1.0]
  ($(zeta.south east)!0.45!(zeta.north east)$)
  node[lab, pos=0.5, yshift=-210, xshift=215, eExt, align=center]
    {$\displaystyle\zeta(s) = -\frac{\Gamma(1-s)}{2\pi i}\oint_{\mathrm{H}}\frac{(-z)^{s-1}}{e^z-1}\,\mathrm{d}z$\\
     解析延拓到 $\mathbb{C}\setminus\{1\}$，Hankel 围道消去分支切割};

\end{scope}% 主图

% ==================== 图例（取自 A-2 详细文案 + 分阶段灰白色） ====================
\node[draw=black!40, fill=gray!8, rounded corners=8pt, thick, align=center, inner sep=10pt, font=\tiny]
  (legend) at (0, -23.5) {
  \begin{tabular}{@{}ll@{\hspace{1.0cm}}ll@{}}
    \multicolumn{2}{l}{\scriptsize\textbf{节点颜色分类图例}} &
    \multicolumn{2}{l}{\scriptsize\textbf{知识谱系连线语义图例}} \\[6pt]
    \tikz[baseline=-0.5ex]\fill[lgPrime, draw=black!70, thick, rounded corners=1.5pt] (0,-0.12) rectangle (0.55,0.12); &
    \textcolor{lgTprime}{\textbf{素数计数函数}（实变，如 $\pi(x)$）} &
    \textcolor{lgEArith}{\rule[0.5ex]{16pt}{1.4pt}} &
    \textcolor{lgTarith}{\textbf{算术互逆}（M\"{o}bius / Stieltjes 反演）} \\[3pt]
    \tikz[baseline=-0.5ex]\fill[lgMu, draw=black!70, thick, rounded corners=1.5pt] (0,-0.12) rectangle (0.55,0.12); &
    \textcolor{lgTmu}{\textbf{算术积性函数}（离散，如 $\mu(n)$）} &
    \textcolor{lgEDef}{\rule[0.5ex]{16pt}{1.4pt}} &
    \textcolor{lgTdef}{\textbf{对象定义}（构造定义与主项提取）} \\[3pt]
    \tikz[baseline=-0.5ex]\fill[lgLi, draw=black!70, thick, rounded corners=1.5pt] (0,-0.12) rectangle (0.55,0.12); &
    \textcolor{lgTli}{\textbf{渐近逼近主项}（如 $\Li(x)$）} &
    \textcolor{lgEElem}{\rule[0.5ex]{16pt}{1.4pt}} &
    \textcolor{lgTelem}{\textbf{初等计数变换}（数论函数基本恒等式）} \\[3pt]
    \tikz[baseline=-0.5ex]\fill[lgXi, draw=black!55, thick, rounded corners=1.5pt] (0,-0.12) rectangle (0.55,0.12); &
    \textcolor{lgTxi}{\textbf{辅助复变函数}（如 $\xi(s)$）} &
    \textcolor{lgEZero}{\rule[0.5ex]{16pt}{1.4pt}} &
    \textcolor{lgTzero}{\textbf{零点结构与猜想}（分布特征及等价命题）} \\[3pt]
    \tikz[baseline=-0.5ex]\fill[lgRH, draw=black!55, thick, rounded corners=1.5pt] (0,-0.12) rectangle (0.55,0.12); &
    \textcolor{lgTrh}{\textbf{核心猜想与目标}（如 RH）} &
    \textcolor{lgEAnal}{\rule[0.5ex]{16pt}{1.4pt}} &
    \textcolor{lgTanal}{\textbf{分析桥梁}（积分变换与 Dirichlet 级数）} \\[3pt]
    \tikz[baseline=-0.5ex]{\clip[rounded corners=1.5pt] (0,-0.12) rectangle (0.55,0.12); \fill[lgGL] (0,-0.12) rectangle (0.55,0); \fill[lgGU] (0,0) rectangle (0.55,0.12); \draw[black!55, thick, rounded corners=1.5pt] (0,-0.12) rectangle (0.55,0.12); \draw[black!55, thick] (0,0) -- (0.55,0);} &
    \textcolor{lgTgamma}{\textbf{解析延拓复变}（初始域/延拓域）} &
    \textcolor{lgEPNT}{\rule[0.5ex]{16pt}{1.4pt}} &
    \textcolor{lgTpnt}{\textbf{渐近定理}（PNT 等确立的渐近推论）} \\[3pt]
    & & \textcolor{lgEExt}{\rule[0.5ex]{16pt}{1.4pt}} &
    \textcolor{lgText}{\textbf{解析延拓}（复围道积分与全纯延拓）} \\[3pt]
    & & \tikz[baseline=-0.5ex]{\draw[lgEZero, line width=1.2pt] (0.00,0) -- (0.12,0); \draw[lgEZero, line width=1.2pt] (0.17,0) -- (0.29,0); \draw[lgEZero, line width=1.2pt] (0.34,0) -- (0.46,0);} &
    \textbf{未证}
  \end{tabular}
};

\end{scope}% 位移：左 1.0cm、上 2.0cm

\end{tikzpicture}

% 几何与 黎曼猜想知识谱系关系图-A-2.tex 同步；未涉及者仅灰白化（含图例）
%
% \spectrumStage 仅为内部进度码（历史命名），不是章号：
%   2  → 第一次着色：第2章 初等算术层（fig:spectrum-stage2）
%   8  → 第二次着色：第9章 延拓/函数方程（fig:spectrum-stage8）
%  16  → 第三次着色：第15章 全文总图（fig:function-relations）
% 正文表述用「第一次 / 第二次 / 第三次着色」或章标签，勿写「第8章图」「第16章图」。
\providecommand{\spectrumStage}{16}
\makeatletter
\colorlet{cDim}{gray!16}\colorlet{dDim}{gray!48}\colorlet{tDim}{gray!52}
\ifnum\spectrumStage=2
  \colorlet{cPrime}{blue!12}\colorlet{dPrime}{black}\colorlet{tPrime}{black}
  \colorlet{cMu}{teal!15}\colorlet{dMu}{teal!60!black}\colorlet{tMu}{black}
  \colorlet{cLi}{yellow!25}\colorlet{dLi}{orange!70!black}\colorlet{tLi}{black}
  \colorlet{cGammaU}{gray!16}\colorlet{cGammaL}{gray!16}\colorlet{dGamma}{gray!48}\colorlet{tGamma}{gray!52}
  \colorlet{cZetaU}{gray!16}\colorlet{cZetaL}{gray!16}\colorlet{tZeta}{gray!52}
  \colorlet{cXi}{gray!16}\colorlet{dXi}{gray!48}\colorlet{tXi}{gray!52}
  \colorlet{cNT}{gray!16}\colorlet{dNT}{gray!48}\colorlet{tNT}{gray!52}
  \colorlet{cProd}{gray!16}\colorlet{dProd}{gray!48}\colorlet{tProd}{gray!52}
  \colorlet{cRH}{gray!16}\colorlet{dRH}{gray!48}\colorlet{tRH}{gray!52}
  \colorlet{eArith}{blue!70!black}\colorlet{eElem}{violet!80!black}\colorlet{eDef}{black!80}
  \colorlet{ePNT}{brown!70!black}\colorlet{eAnal}{gray!42}\colorlet{eExpl}{gray!42}
  \colorlet{eExt}{gray!42}\colorlet{eZero}{gray!42}\colorlet{eLate}{gray!42}
  \colorlet{lgPrime}{blue!12}\colorlet{lgMu}{teal!15}\colorlet{lgLi}{yellow!25}
  \colorlet{lgXi}{gray!22}\colorlet{lgRH}{gray!22}\colorlet{lgGU}{gray!22}\colorlet{lgGL}{gray!22}
  \colorlet{lgEArith}{blue!70!black}\colorlet{lgEDef}{black!80}\colorlet{lgEElem}{violet!80!black}
  \colorlet{lgEZero}{gray!45}\colorlet{lgEAnal}{gray!45}\colorlet{lgEPNT}{brown!70!black}\colorlet{lgEExt}{gray!45}
  \colorlet{lgTprime}{black}\colorlet{lgTmu}{black}\colorlet{lgTli}{black}
  \colorlet{lgTxi}{gray!50}\colorlet{lgTrh}{gray!50}\colorlet{lgTgamma}{gray!50}
  \colorlet{lgTarith}{black}\colorlet{lgTdef}{black}\colorlet{lgTelem}{black}
  \colorlet{lgTzero}{gray!50}\colorlet{lgTanal}{gray!50}\colorlet{lgTpnt}{black}\colorlet{lgText}{gray!50}
\else
\ifnum\spectrumStage=8
  \colorlet{cPrime}{blue!12}\colorlet{dPrime}{black}\colorlet{tPrime}{black}
  \colorlet{cMu}{teal!15}\colorlet{dMu}{teal!60!black}\colorlet{tMu}{black}
  \colorlet{cLi}{yellow!25}\colorlet{dLi}{orange!70!black}\colorlet{tLi}{black}
  \colorlet{cGammaU}{blue!15}\colorlet{cGammaL}{red!15}\colorlet{dGamma}{black}\colorlet{tGamma}{black}
  \colorlet{cZetaU}{blue!15}\colorlet{cZetaL}{red!15}\colorlet{tZeta}{black}
  \colorlet{cXi}{gray!16}\colorlet{dXi}{gray!48}\colorlet{tXi}{gray!52}
  \colorlet{cNT}{gray!16}\colorlet{dNT}{gray!48}\colorlet{tNT}{gray!52}
  \colorlet{cProd}{gray!16}\colorlet{dProd}{gray!48}\colorlet{tProd}{gray!52}
  \colorlet{cRH}{gray!16}\colorlet{dRH}{gray!48}\colorlet{tRH}{gray!52}
  \colorlet{eArith}{blue!70!black}\colorlet{eElem}{violet!80!black}\colorlet{eDef}{black!80}
  \colorlet{ePNT}{brown!70!black}\colorlet{eAnal}{green!50!black}\colorlet{eExpl}{gray!42}
  \colorlet{eExt}{red!70!black}\colorlet{eZero}{gray!42}\colorlet{eLate}{gray!42}
  \colorlet{lgPrime}{blue!12}\colorlet{lgMu}{teal!15}\colorlet{lgLi}{yellow!25}
  \colorlet{lgXi}{gray!22}\colorlet{lgRH}{gray!22}\colorlet{lgGU}{blue!15}\colorlet{lgGL}{red!15}
  \colorlet{lgEArith}{blue!70!black}\colorlet{lgEDef}{black!80}\colorlet{lgEElem}{violet!80!black}
  \colorlet{lgEZero}{gray!45}\colorlet{lgEAnal}{green!50!black}\colorlet{lgEPNT}{brown!70!black}\colorlet{lgEExt}{red!70!black}
  \colorlet{lgTprime}{black}\colorlet{lgTmu}{black}\colorlet{lgTli}{black}
  \colorlet{lgTxi}{gray!50}\colorlet{lgTrh}{gray!50}\colorlet{lgTgamma}{black}
  \colorlet{lgTarith}{black}\colorlet{lgTdef}{black}\colorlet{lgTelem}{black}
  \colorlet{lgTzero}{gray!50}\colorlet{lgTanal}{black}\colorlet{lgTpnt}{black}\colorlet{lgText}{black}
\else
  \colorlet{cPrime}{blue!12}\colorlet{dPrime}{black}\colorlet{tPrime}{black}
  \colorlet{cMu}{teal!15}\colorlet{dMu}{teal!60!black}\colorlet{tMu}{black}
  \colorlet{cLi}{yellow!25}\colorlet{dLi}{orange!70!black}\colorlet{tLi}{black}
  \colorlet{cGammaU}{blue!15}\colorlet{cGammaL}{red!15}\colorlet{dGamma}{black}\colorlet{tGamma}{black}
  \colorlet{cZetaU}{blue!15}\colorlet{cZetaL}{red!15}\colorlet{tZeta}{black}
  \colorlet{cXi}{cyan!15}\colorlet{dXi}{cyan!60!black}\colorlet{tXi}{black}
  \colorlet{cNT}{cyan!15}\colorlet{dNT}{cyan!60!black}\colorlet{tNT}{black}
  \colorlet{cProd}{cyan!15}\colorlet{dProd}{cyan!60!black}\colorlet{tProd}{black}
  \colorlet{cRH}{orange!15}\colorlet{dRH}{orange!70!black}\colorlet{tRH}{black}
  \colorlet{eArith}{blue!70!black}\colorlet{eElem}{violet!80!black}\colorlet{eDef}{black!80}
  \colorlet{ePNT}{brown!70!black}\colorlet{eAnal}{green!50!black}\colorlet{eExpl}{green!50!black}
  \colorlet{eExt}{red!70!black}\colorlet{eZero}{orange!80!black}\colorlet{eLate}{black!80}
  \colorlet{lgPrime}{blue!12}\colorlet{lgMu}{teal!15}\colorlet{lgLi}{yellow!25}
  \colorlet{lgXi}{cyan!15}\colorlet{lgRH}{orange!15}\colorlet{lgGU}{blue!15}\colorlet{lgGL}{red!15}
  \colorlet{lgEArith}{blue!70!black}\colorlet{lgEDef}{black!80}\colorlet{lgEElem}{violet!80!black}
  \colorlet{lgEZero}{orange!80!black}\colorlet{lgEAnal}{green!50!black}\colorlet{lgEPNT}{brown!70!black}\colorlet{lgEExt}{red!70!black}
  \colorlet{lgTprime}{black}\colorlet{lgTmu}{black}\colorlet{lgTli}{black}
  \colorlet{lgTxi}{black}\colorlet{lgTrh}{black}\colorlet{lgTgamma}{black}
  \colorlet{lgTarith}{black}\colorlet{lgTdef}{black}\colorlet{lgTelem}{black}
  \colorlet{lgTzero}{black}\colorlet{lgTanal}{black}\colorlet{lgTpnt}{black}\colorlet{lgText}{black}
\fi\fi
\makeatother

\begin{tikzpicture}[scale=1.0, transform shape,
  box/.style    = {draw, rounded corners=5pt, minimum width=2.8cm,
                   minimum height=0.9cm, align=center, font=\small, thick},
  zbox/.style   = {draw, rounded corners=5pt, minimum width=2.8cm,
                   minimum height=0.9cm, align=center, font=\small,
                   thick, fill=gray!18},
  splitbox/.style = {draw, rounded corners=5pt, minimum width=7cm,
                     minimum height=2.2cm, align=center, font=\small, thick,
                     rectangle split, rectangle split parts=2,
                     rectangle split part fill={red!8, white}},
  arr/.style    = {-{Stealth[length=6pt,width=5pt]}, thick},
  barr/.style   = {{Stealth[length=6pt,width=5pt]}-{Stealth[length=6pt,width=5pt]}, thick},
  % 线型：虚线仅用于 N(T)→RH、ξ↔RH；其余为实线
  proven/.style = {solid},
  cond/.style   = {dashed},
  lab/.style    = {font=\tiny, align=center, inner sep=3pt},
  rhbox/.style  = {draw=dRH, rounded corners=5pt, minimum width=8cm,
                   minimum height=1.0cm, align=center, font=\small,
                   fill=cRH, text=tRH, thick}
]

% 布局与 A-2 同步：左移 1.0cm、上移 2.0cm
% 左右略放宽，避免 Γ/ζ 自环旁公式标注被裁切
\path[use as bounding box]
  (-10.2, -26.2) rectangle (10.0, 7.4);

\begin{scope}[shift={(-1.0,2.0)}]

\node[font=\bfseries\Large, align=center] at (0, 4.35)
  {黎曼猜想知识谱系关系图};

% 主图单独上移 1cm；标题与图例位置不动
\begin{scope}[shift={(-2.75,-1.2)}]
%% ==================== 节点分层 ====================

% 第一层：素数计数函数 (y=0)
\node[box, fill=cPrime, draw=dPrime]    (J)    at (-4.2,  0.0) {$J(x)$};
\node[box, fill=cPrime, draw=dPrime]    (pi)   at ( 0.0,  0.0) {$\pi(x)$};
\node[box, fill=cPrime, draw=dPrime]    (th)   at ( 5.0,  0.0) {$\theta(x)$};
\node[box, fill=cPrime, draw=dPrime]    (ps)   at ( 9.7,  0.0) {$\psi(x)$};

% 第二层：积性函数 (y=-3.0)
% 节点定义修改
\node[box, fill=cMu, draw=dMu] (mu)  at (0.3, -3.0) {$\mu(n)$};
\node[box, fill=cMu, draw=dMu] (Lam) at (5.0, -3.0) {$\Lambda(n)$};

%Weierstrass乘积
\node[box, fill=cProd, draw=dProd, text=tProd]  (prod)   at ( 9.5,  -17.0) {\tiny\text{Hadamard 乘积}};
% M(x) 节点已移除：初等等价信息直接标注在 RH 框内

% 第三层：Li 与 Gamma (y=-6.0 至 -7.0)
\node[box, fill=cLi, draw=dLi] (Li) at (-2.0, -6.0) {$\Li(x)$};

% Gamma(s) 上下分色：上部标收敛域，下部标延拓
\node[splitbox, rectangle split draw splits=false, minimum width=5cm,
      minimum height=4.7cm,
      rectangle split part fill={cGammaU, cGammaL}] (gamma) at (7.0, -6.0)
    {
     \centering \color{tGamma}$\operatorname{Re}(s)>0$
     \nodepart{two}
     \raisebox{6pt}{\color{tGamma}$\mathbb{C}\setminus\{0,-1,-2,\dots\}$}
    };
% Gamma(s) 标注（固定在节点中线左侧）
\node[font=\small, left=-30pt, text=tGamma] at (gamma.west) {$\Gamma(s)$};

% 【修订3】ζ(s) 节点：下半部分补全函数方程左侧 ζ(s)=
\node[splitbox, rectangle split draw splits=false, minimum height=4cm,
      rectangle split part fill={cZetaU, cZetaL}] (zeta) at (0.3, -10.5)
    {\centering\small
     \color{tZeta}$\displaystyle
       \sum_{n=1}^{\infty}\frac{1}{n^s}
       =\prod_{p}\dfrac{1}{1-p^{-s}}
       \quad(\operatorname{Re}(s)>1)$%
     \vphantom{$\displaystyle\sum_{n=1}^{\infty}$}%
     \nodepart{two}
     % 修订：补全 ζ(s)= 使函数方程完整
     \textcolor{tZeta}{$\displaystyle\zeta(s)=
       2^s\pi^{s-1}
       \sin\!\left(\frac{\pi s}{2}\right)\Gamma(1-s)\,\zeta(1-s)$%
     \vphantom{$\displaystyle\sum_{n=1}^{\infty}\mathbb{C}\setminus\{1\}$}}
    };
% ζ(s) 标注


% ξ(s) 节点：下移至与 RH 同一水平，右侧对称于 M(x)
\node[box, fill=cXi, draw=dXi, text=tXi, minimum height=1.2cm]
  (xi) at (4.85, -17.0)
  {$\xi(s)$\\ \tiny 非平凡零点$\rho$与$\zeta(s)$相同\\
  \tiny $\xi(s)=\xi(1-s)$(函数方程推论)
  };


% 【修订8】N(T) 节点：置于 ζ 与 RH 正中间，构成竖直逻辑链
\node[box, fill=cNT, draw=dNT, text=tNT, minimum height=1.3cm]
  (NT) at (2.3, -14.0)
  {$N(T)$\\ \tiny $\zeta(s)$ 零点计数函数\\[1pt]\tiny（幅角原理）};

% ==================== 黎曼猜想节点 ====================
% RH 下移至 y=-20.5，为 N(T) 留出充分空间
\node[rhbox, minimum width=4.6cm, minimum height=2.2cm] (rh) at (-1.9, -17.0) {
    {\tiny\color{tRH} 黎曼猜想 (Riemann Hypothesis, 1859)}\\[3pt]
    {\tiny\color{tRH}$\zeta(s)$ 的所有非平凡零点 $\rho$ 满足
    $\operatorname{Re}(\rho)=\dfrac{1}{2}$}\\[4pt]
{\tiny\color{tRH} 初等等价（梅滕斯函数）：%
  $\displaystyle M(x)=\sum_{n\leq x}\mu(n)$}\\
{\tiny\color{tRH}
  $\mathrm{RH} \iff M(x) = O(x^{1/2+\varepsilon}),\ \forall \varepsilon>0$}
};


% ==================== 连线 ====================

% ----- 第一层内部 -----
\draw[barr, eArith] (J.east) -- (pi.west)
  node[lab, midway, above=14pt, xshift=-2pt, eArith]
    {$\displaystyle\pi(x)=\sum_{n=1}^{\infty}\frac{\mu(n)}{n}J(x^{1/n})$}
  node[lab, midway, below=14pt, xshift=-2pt, eArith]
    {$\displaystyle J(x)=\sum_{n=1}^{\infty}\frac{1}{n}\pi(x^{1/n})$};

\draw[barr, eElem] (pi.east) -- (th.west)
  node[lab, midway, above=14pt, eElem]
    {$\displaystyle\theta(x)=\pi(x)\log x-\int_2^x\frac{\pi(t)}{t}\,\mathrm{d}t$}
  node[lab, midway, above={14pt}, yshift={-50pt}, eElem]
    {\tiny $\displaystyle\pi(x) = \dfrac{\theta(x)}{\log x}+ \int_{2}^{x} \frac{\theta(t)}{t\,(\log t)^2} \,\mathrm{d}t$};

\draw[barr, eElem] (th.east) -- (ps.west)
  node[lab, midway, above=11pt, xshift={5pt}, eElem]
    {$\displaystyle\psi(x)=\sum_{k=1}^{\infty}\theta(x^{1/k})$}
  node[lab, midway, above=-31pt, xshift={5pt}, yshift={-8pt}, eElem]
    {\tiny $\displaystyle\theta(x)=\sum_{k=1}^{\infty}\mu(k)\psi(x^{1/k})$};


% ----- 第二层内部 -----
\draw[barr, eArith] (mu.east) -- (Lam.west)
  node[lab, midway, yshift=20, xshift=0.3cm, eArith]
    {$\displaystyle\Lambda(n)=-\sum_{d\mid n}\mu(d)\log d$
     \quad{\rm (M\"{o}bius 反演)}}
  node[lab, midway, yshift={-12pt-0.5cm}, xshift=0, eArith]
    {$\displaystyle\log n=\sum_{d\mid n}\Lambda(d)$};

% μ(n) → M(x) 的定义连线已省去（见节点内文字），避免与已有连线交叉

% Lambda -> psi
\draw[arr, eDef] (Lam.east) to[out=0, in=-90, looseness=1.0]
  ($(ps.south west)+(30pt,0)$)
  node[lab, pos=0.5, yshift=-90, xshift=250, eDef]
    {$\displaystyle\psi(x)=\sum_{n\le x}\Lambda(n)$};

% ----- 第三层：Li 和 Gamma -----

% 【修订5】Li(x) → π(x)：改为渐近等价语义（非因果推导）
\draw[arr, ePNT]
  ($(Li.north)!.75!(Li.north east)$)
  to[out=90, in=-135, looseness=1.8]
  ($(pi.south west)+(35pt,0)$)
  node[lab, pos=0.4, yshift=-45, xshift=45, ePNT, align=left]
    {\textbf{素数定理（PNT）}：$\pi(x)\sim\Li(x)$\\
     (若RH成立，$\pi(x)=\Li(x)+O(\sqrt{x}\log x)$)};

% Li -> J（显式公式主项；属 eExpl：第2/9章灰白，第15章全文总图着色）
\draw[arr, eExpl]
  (Li.north) to[out=135, in=-90, looseness=1.5]
  ($(J.south east)+(-22pt,0)$)
  node[lab, pos=0.5, yshift=-115, xshift=-65, eExpl, align=center]
    {$\Li(x)$ 是 $J(x)$ 显式\\[2pt]公式的主项};

% Gamma（Re>0）→ ζ：Euler 积分表示（该积分公式本身要求 Re(s)>1）
\draw[arr, eAnal]
  ([yshift=6pt] gamma.west)
  to[out=180, in=90, looseness=0.8]
  ($(zeta.north)+(30pt,0)$)
  node[lab, pos=0.5, yshift=-235, xshift=90,eAnal, align=center]
    {$\displaystyle\zeta(s)\Gamma(s)
      =\int_0^{\infty}\frac{x^{s-1}}{e^x-1}\,\mathrm{d}x$\\
     $(\operatorname{Re}(s)>1)$};

% μ(n) → ζ(s)
\draw[arr, eAnal]
  (mu.south) to[out=-90, in=90, looseness=1.2]
  ($(zeta.north)+({-20pt+0.3cm},0)$)
  node[lab, pos=0.5, yshift={-120pt-0.5cm-0.4cm}, xshift={45pt-5pt}, eAnal]
    {$\displaystyle\frac{1}{\zeta(s)}=\sum_{n=1}^{\infty}\frac{\mu(n)}{n^s}$\\
     $(\operatorname{Re}(s)>1)$};

% ----- J(x) 与 ψ(x) 的 Stieltjes 积分互逆关系 -----
% 【修订4】加注 dψ(t)/ln t 在 t=1 附近需 t>1 的说明
\draw[barr, eArith, line width=0.8pt, rounded corners=8pt]
  ($(J.north)+(0,0.0cm)$)
  -- ++(0,2.0cm)
  -- ++(13.5cm,0)
  -- ++(0,-2.0cm)
  node[lab, pos=0.35, above=6pt, yshift={10pt+0.5cm}, xshift=-200, eArith, align=center]
    {$\displaystyle\psi(x)=\int_{2^-}^x\log t\,\mathrm{d}J(t)$\quad(Stieltjes 积分，$J\to\psi$)}
  node[lab, pos=0.65, below=6pt, yshift={32pt+0.2cm}, xshift=-205, eArith, align=center]
    {$\displaystyle J(x)=\int_{2^-}^x\frac{\mathrm{d}\psi(t)}{\log t}$
     \quad(Stieltjes 反演，$\psi\to J$；积分自 $t>1$)};

% ----- 第四层：ζ 和 ξ 与第二层的连线 -----

% ζ → ψ：显式公式（Perron 积分 + 留数，单向；x≠p^k 时公式精确成立）
\draw[arr, eExpl, rounded corners=30pt]
  ($(zeta.east)+(0, -10pt)$) -| ($(ps.south)+(30pt,0)$)
  node[lab, pos=0.85, right=10pt, yshift=-170, xshift={-170pt-0.2cm},eExpl, align=left]
{$\psi(x)$ 显式公式（Perron 积分，$\zeta\to\psi$）\\
     $\displaystyle \psi(x) = \frac{1}{2\pi i}\int_{c-i\infty}^{c+i\infty}\!\left(-\frac{\zeta'(s)}{\zeta(s)}\right)\!\frac{x^s}{s}\,\mathrm{d}s$\\
     $\displaystyle = x-\sum_{\rho}\frac{x^{\rho}}{\rho}-\log(2\pi)-\tfrac{1}{2}\log(1-x^{-2})$\\
     \tiny $x>1,\;x\neq p^k$；$\rho$ 按 $|\mathrm{Im}(\rho)|$ 升序配对，条件收敛};

% ψ → ζ：Mellin 变换（单向）
\draw[arr, eExpl, rounded corners=30pt]
  ($(ps.south)+(40pt,0)$) |- ($(zeta.east)+(0, -20pt)$)
  node[lab, pos=0.8, yshift=-27, xshift=30,eExpl, align=center]
    {Mellin 变换（$\psi(x) \longleftrightarrow -\dfrac{\zeta'(s)}{\zeta(s)}$，即负对数导数）\\
     $\displaystyle -\frac{\zeta'(s)}{\zeta(s)} = s\int_{1}^{\infty}\frac{\psi(x)}{x^{s+1}}\,\mathrm{d}x$};
       
       
       
       
       
       
       
% Lambda -> zeta（对数微分 Dirichlet 级数）
\draw[arr, eAnal]
  ($(Lam.south west)+(6pt,0)$) to[out=-90, in=90, looseness=0.4]
  ($(zeta.north west)!0.50!(zeta.north east)$)
  node[lab, pos=0.5, yshift={-120pt-0.5cm-0.4cm+1cm}, xshift={145pt-5pt}, eAnal]
    {$\displaystyle-\frac{\zeta'(s)}{\zeta(s)}
      =\sum_{n=1}^{\infty}\frac{\Lambda(n)}{n^s}$};

% J(x) 显式公式（ζ → J）
\draw[arr, eExpl, rounded corners=50pt]
  ($(zeta.south west)+(0,6pt)$) -- ++(-60pt,0) -- ($(J.south west)+(5pt,0)$)
  node[lab, pos=0.5, yshift=-203, xshift={75pt-0.2cm}, eExpl, align=center]
{\textbf{Perron 公式} $\Rightarrow$ $J(x)$ 显式展开\\
     $\displaystyle
     \begin{aligned}
         &J(x)=\frac{1}{2\pi i}\int_{c-i\infty}^{c+i\infty}
           \frac{\log\zeta(s)}{s}\,x^s\,\mathrm{d}s\\
         &=\Li(x)-\sum_{\rho}\Li(x^{\rho})
           -\log 2+\int_x^{\infty}\frac{\mathrm{d}t}{t(t^2-1)\log t}
     \end{aligned}$\\
     \tiny $x>1,\;x\neq p^k$；$\rho$ 按 $|\mathrm{Im}(\rho)|$ 升序配对，条件收敛};

% ζ → J：Mellin 变换（ln ζ）
\draw[arr, eAnal, rounded corners=5pt]
  (zeta.north west) to[out=110, in=-90, looseness=1.5] (J.south)
  node[lab, pos=0.4, yshift=-205, xshift=-55, eAnal, align=center]
    {$\displaystyle\log\zeta(s)=s\int_1^{\infty}J(x)\,x^{-s-1}\,\mathrm{d}x$\\
     $(\operatorname{Re}(s)>1)$};

% ζ → ξ（定义）：ξ 已下移，连线沿右侧空白区域向右下延伸
\draw[arr, eLate, rounded corners=15pt]
  ($(zeta.south east)!0.15!(zeta.north east)$)
  -- ++(1.0, 0)
  -- ++(0, -3.8)
  -- ($(xi.north)+(0pt,0)$)
  node[lab, pos=0.25, right=4pt,yshift=30, xshift=0, eLate, align=center]
    {$\xi(s)=\dfrac{1}{2}s(s-1)\pi^{-s/2}\Gamma(s/2)\zeta(s)$  
};






% xi -> prod (Weierstrass product connection)
\draw[arr, eLate]
  %(xi.south) -- (prod.south)
  (xi.east) --(prod.west)
  node[midway, right=4pt, yshift=28, xshift=-45, eLate]
    {\tiny$\displaystyle\xi(s)=e^{A+Bs}\prod_{\rho}\left(1-\frac{s}{\rho}\right)e^{s/\rho}$};

% prod -> ps (zero sum in explicit formula)
%\draw[arr, violet!70!black, rounded corners=15pt]
  %(prod.east) -- ++(4.4, 0) |- (ps.south)
  %node[lab, pos=0.6, right=2pt, violet!70!black, align=left]
   % {对数微分与 Perron 积分\\
    % 给出显式公式中\\
    % 零点和 $\sum_{\rho}\frac{x_{\rho}}{\rho}$ 项};
%
% ζ 自身左侧弧线（解析延拓）：与表 tab:spectrum-stage8-guide 一致
\draw[arr, eExt, rounded corners=5pt]
  ($(zeta.north west)!0.15!(zeta.south west)$)
  to[out=180, in=180, looseness=1.6]
  ($(zeta.north west)!0.85!(zeta.south west)$);
% 标签用独立节点锚定在 ζ 西侧弧线旁（避免 path 末 node 漂移）
% 单行；中心在 zeta.west 右 2.5cm（相对初版再累计右移）
\node[lab, eExt, align=center, anchor=center]
  at ([xshift=2.5cm]zeta.west)
  {$\zeta$ 亚纯延拓至 $\mathbb{C}\setminus\{1\}$（Hankel / $\vartheta$ 两法）};

% ζ → N(T)：竖直向下（弧线）
\draw[arr, eZero]
  ($(zeta.south east)!0.10!(zeta.north east)$)
  to[out=0, in=90] ++(0.8, -0.8)
  to[out=-90, in=0] (NT.east)
  node[lab, pos=0.6, right=6pt, yshift=-357, xshift=5, eZero, align=center]
    {$\displaystyle N(T)=\frac{T}{2\pi}\log\frac{T}{2\pi e}+O(\log T)$\\
     \tiny 零点计数公式（von Mangoldt）};

% 【已删去】ζ → RH 蓝色直连线（设计意图改为ζ→N(T)→RH线性链）
%\draw[arr, eArith, rounded corners=10pt]
%  (zeta.south)
%  -- ++(0, -3.35)
%  -- ++(-4.1, 0)
%  -- ($(rh.north)!0.80!(rh.north west)$)
%  node[lab, pos=0.35, left=6pt, yshift=20pt,xshift=120,eArith, align=center]
%    {
%     $\zeta(s)$ 的非平凡零点全部位于直线 $\operatorname{Re}(s)=\dfrac12$ 上};
% RH -- P(s) 画线 + 文字
%\draw[arr, cyan!60!black]
 % ($(rh.east)!0.60!(rh.south east)$) -- (prod.west)    % 起点 -- 终点（必须有！）
  %node[midway, right=6pt, font=\tiny, cyan!60!black, align=center]
  %  {$\zeta(s) = \frac{e^{s(\log(2\pi) - 1)}}{2(s-1)} \prod_{n=1}^{\infty} \left(1 + \frac{s}{2n}\right)e^{-s/2n} P(s)$};

% ξ(s) ↔ RH、N(T) → RH：虚线标注（与 RH 相关的谱系连线）

% N(T) → RH：完成 ζ→N(T)→RH 线性逻辑链（虚线）
\draw[arr, eZero, cond, rounded corners=8pt]
  (NT.south west)
  to[out=-150, in=60, looseness=0.9]
  ($(rh.north)+(20pt, 0)$)
node[lab, pos=0.55, left=4pt, yshift=-430, xshift=-25, eZero, align=center]
    {$\mathrm{RH} \iff N_0(T)=N(T),\ \forall\,T>0$\\
   \tiny（$N_0(T)$：临界线 $\operatorname{Re}(s)=\dfrac{1}{2}$ 上虚部 $\leq T$ 的零点数）};

\draw[barr, eZero, cond]
  (xi.west) -- ($(rh.east)!0.0!(rh.south east)$)
  node[lab, midway, above=20pt,yshift=3pt,xshift=5, eZero, align=center]
    {$\mathrm{RH}\iff\xi\!\left(\dfrac{1}{2}+it\right)$\\
     的所有零点 $t$ 均为实数};

% M(x)→RH 箭头已删去：等价信息直接写在 RH 框内

% Gamma 自身解析延拓弧线：与表 tab:spectrum-stage8-guide 一致
\draw[arr, eExt, rounded corners=1.5pt]
  ($(gamma.north east)+(0,-5pt)$)
  to[out=0, in=0, looseness=1.6]
  ($(gamma.south east)+(0,5pt)$);
% 标签用独立节点锚定在 Γ 东侧弧线旁（仅递推式；延拓域已在框下半）
% 相对初版 (gamma.east+14pt)：净向左 3.7cm，上移 2pt
\node[lab, eExt, align=center, anchor=west]
  at ($(gamma.east)+({14pt-3.7cm},2pt)$)
  {递推延拓 $\Gamma(s)=\Gamma(s+1)/s$};

% Gamma（延拓后）→ ζ：Hankel 围道
\draw[arr, eExt]
  ($(gamma.south west)+(40pt,0)$)
  to[out=180, in=0, looseness=1.0]
  ($(zeta.south east)!0.45!(zeta.north east)$)
  node[lab, pos=0.5, yshift=-210, xshift=215, eExt, align=center]
    {$\displaystyle\zeta(s) = -\frac{\Gamma(1-s)}{2\pi i}\oint_{\mathrm{H}}\frac{(-z)^{s-1}}{e^z-1}\,\mathrm{d}z$\\
     解析延拓到 $\mathbb{C}\setminus\{1\}$，Hankel 围道消去分支切割};

\end{scope}% 主图

% ==================== 图例（取自 A-2 详细文案 + 分阶段灰白色） ====================
\node[draw=black!40, fill=gray!8, rounded corners=8pt, thick, align=center, inner sep=10pt, font=\tiny]
  (legend) at (0, -23.5) {
  \begin{tabular}{@{}ll@{\hspace{1.0cm}}ll@{}}
    \multicolumn{2}{l}{\scriptsize\textbf{节点颜色分类图例}} &
    \multicolumn{2}{l}{\scriptsize\textbf{知识谱系连线语义图例}} \\[6pt]
    \tikz[baseline=-0.5ex]\fill[lgPrime, draw=black!70, thick, rounded corners=1.5pt] (0,-0.12) rectangle (0.55,0.12); &
    \textcolor{lgTprime}{\textbf{素数计数函数}（实变，如 $\pi(x)$）} &
    \textcolor{lgEArith}{\rule[0.5ex]{16pt}{1.4pt}} &
    \textcolor{lgTarith}{\textbf{算术互逆}（M\"{o}bius / Stieltjes 反演）} \\[3pt]
    \tikz[baseline=-0.5ex]\fill[lgMu, draw=black!70, thick, rounded corners=1.5pt] (0,-0.12) rectangle (0.55,0.12); &
    \textcolor{lgTmu}{\textbf{算术积性函数}（离散，如 $\mu(n)$）} &
    \textcolor{lgEDef}{\rule[0.5ex]{16pt}{1.4pt}} &
    \textcolor{lgTdef}{\textbf{对象定义}（构造定义与主项提取）} \\[3pt]
    \tikz[baseline=-0.5ex]\fill[lgLi, draw=black!70, thick, rounded corners=1.5pt] (0,-0.12) rectangle (0.55,0.12); &
    \textcolor{lgTli}{\textbf{渐近逼近主项}（如 $\Li(x)$）} &
    \textcolor{lgEElem}{\rule[0.5ex]{16pt}{1.4pt}} &
    \textcolor{lgTelem}{\textbf{初等计数变换}（数论函数基本恒等式）} \\[3pt]
    \tikz[baseline=-0.5ex]\fill[lgXi, draw=black!55, thick, rounded corners=1.5pt] (0,-0.12) rectangle (0.55,0.12); &
    \textcolor{lgTxi}{\textbf{辅助复变函数}（如 $\xi(s)$）} &
    \textcolor{lgEZero}{\rule[0.5ex]{16pt}{1.4pt}} &
    \textcolor{lgTzero}{\textbf{零点结构与猜想}（分布特征及等价命题）} \\[3pt]
    \tikz[baseline=-0.5ex]\fill[lgRH, draw=black!55, thick, rounded corners=1.5pt] (0,-0.12) rectangle (0.55,0.12); &
    \textcolor{lgTrh}{\textbf{核心猜想与目标}（如 RH）} &
    \textcolor{lgEAnal}{\rule[0.5ex]{16pt}{1.4pt}} &
    \textcolor{lgTanal}{\textbf{分析桥梁}（积分变换与 Dirichlet 级数）} \\[3pt]
    \tikz[baseline=-0.5ex]{\clip[rounded corners=1.5pt] (0,-0.12) rectangle (0.55,0.12); \fill[lgGL] (0,-0.12) rectangle (0.55,0); \fill[lgGU] (0,0) rectangle (0.55,0.12); \draw[black!55, thick, rounded corners=1.5pt] (0,-0.12) rectangle (0.55,0.12); \draw[black!55, thick] (0,0) -- (0.55,0);} &
    \textcolor{lgTgamma}{\textbf{解析延拓复变}（初始域/延拓域）} &
    \textcolor{lgEPNT}{\rule[0.5ex]{16pt}{1.4pt}} &
    \textcolor{lgTpnt}{\textbf{渐近定理}（PNT 等确立的渐近推论）} \\[3pt]
    & & \textcolor{lgEExt}{\rule[0.5ex]{16pt}{1.4pt}} &
    \textcolor{lgText}{\textbf{解析延拓}（复围道积分与全纯延拓）} \\[3pt]
    & & \tikz[baseline=-0.5ex]{\draw[lgEZero, line width=1.2pt] (0.00,0) -- (0.12,0); \draw[lgEZero, line width=1.2pt] (0.17,0) -- (0.29,0); \draw[lgEZero, line width=1.2pt] (0.34,0) -- (0.46,0);} &
    \textbf{未证}
  \end{tabular}
};

\end{scope}% 位移：左 1.0cm、上 2.0cm

\end{tikzpicture}

% 几何与 黎曼猜想知识谱系关系图-A-2.tex 同步；未涉及者仅灰白化（含图例）
%
% \spectrumStage 仅为内部进度码（历史命名），不是章号：
%   2  → 第一次着色：第2章 初等算术层（fig:spectrum-stage2）
%   8  → 第二次着色：第9章 延拓/函数方程（fig:spectrum-stage8）
%  16  → 第三次着色：第15章 全文总图（fig:function-relations）
% 正文表述用「第一次 / 第二次 / 第三次着色」或章标签，勿写「第8章图」「第16章图」。
\providecommand{\spectrumStage}{16}
\makeatletter
\colorlet{cDim}{gray!16}\colorlet{dDim}{gray!48}\colorlet{tDim}{gray!52}
\ifnum\spectrumStage=2
  \colorlet{cPrime}{blue!12}\colorlet{dPrime}{black}\colorlet{tPrime}{black}
  \colorlet{cMu}{teal!15}\colorlet{dMu}{teal!60!black}\colorlet{tMu}{black}
  \colorlet{cLi}{yellow!25}\colorlet{dLi}{orange!70!black}\colorlet{tLi}{black}
  \colorlet{cGammaU}{gray!16}\colorlet{cGammaL}{gray!16}\colorlet{dGamma}{gray!48}\colorlet{tGamma}{gray!52}
  \colorlet{cZetaU}{gray!16}\colorlet{cZetaL}{gray!16}\colorlet{tZeta}{gray!52}
  \colorlet{cXi}{gray!16}\colorlet{dXi}{gray!48}\colorlet{tXi}{gray!52}
  \colorlet{cNT}{gray!16}\colorlet{dNT}{gray!48}\colorlet{tNT}{gray!52}
  \colorlet{cProd}{gray!16}\colorlet{dProd}{gray!48}\colorlet{tProd}{gray!52}
  \colorlet{cRH}{gray!16}\colorlet{dRH}{gray!48}\colorlet{tRH}{gray!52}
  \colorlet{eArith}{blue!70!black}\colorlet{eElem}{violet!80!black}\colorlet{eDef}{black!80}
  \colorlet{ePNT}{brown!70!black}\colorlet{eAnal}{gray!42}\colorlet{eExpl}{gray!42}
  \colorlet{eExt}{gray!42}\colorlet{eZero}{gray!42}\colorlet{eLate}{gray!42}
  \colorlet{lgPrime}{blue!12}\colorlet{lgMu}{teal!15}\colorlet{lgLi}{yellow!25}
  \colorlet{lgXi}{gray!22}\colorlet{lgRH}{gray!22}\colorlet{lgGU}{gray!22}\colorlet{lgGL}{gray!22}
  \colorlet{lgEArith}{blue!70!black}\colorlet{lgEDef}{black!80}\colorlet{lgEElem}{violet!80!black}
  \colorlet{lgEZero}{gray!45}\colorlet{lgEAnal}{gray!45}\colorlet{lgEPNT}{brown!70!black}\colorlet{lgEExt}{gray!45}
  \colorlet{lgTprime}{black}\colorlet{lgTmu}{black}\colorlet{lgTli}{black}
  \colorlet{lgTxi}{gray!50}\colorlet{lgTrh}{gray!50}\colorlet{lgTgamma}{gray!50}
  \colorlet{lgTarith}{black}\colorlet{lgTdef}{black}\colorlet{lgTelem}{black}
  \colorlet{lgTzero}{gray!50}\colorlet{lgTanal}{gray!50}\colorlet{lgTpnt}{black}\colorlet{lgText}{gray!50}
\else
\ifnum\spectrumStage=8
  \colorlet{cPrime}{blue!12}\colorlet{dPrime}{black}\colorlet{tPrime}{black}
  \colorlet{cMu}{teal!15}\colorlet{dMu}{teal!60!black}\colorlet{tMu}{black}
  \colorlet{cLi}{yellow!25}\colorlet{dLi}{orange!70!black}\colorlet{tLi}{black}
  \colorlet{cGammaU}{blue!15}\colorlet{cGammaL}{red!15}\colorlet{dGamma}{black}\colorlet{tGamma}{black}
  \colorlet{cZetaU}{blue!15}\colorlet{cZetaL}{red!15}\colorlet{tZeta}{black}
  \colorlet{cXi}{gray!16}\colorlet{dXi}{gray!48}\colorlet{tXi}{gray!52}
  \colorlet{cNT}{gray!16}\colorlet{dNT}{gray!48}\colorlet{tNT}{gray!52}
  \colorlet{cProd}{gray!16}\colorlet{dProd}{gray!48}\colorlet{tProd}{gray!52}
  \colorlet{cRH}{gray!16}\colorlet{dRH}{gray!48}\colorlet{tRH}{gray!52}
  \colorlet{eArith}{blue!70!black}\colorlet{eElem}{violet!80!black}\colorlet{eDef}{black!80}
  \colorlet{ePNT}{brown!70!black}\colorlet{eAnal}{green!50!black}\colorlet{eExpl}{gray!42}
  \colorlet{eExt}{red!70!black}\colorlet{eZero}{gray!42}\colorlet{eLate}{gray!42}
  \colorlet{lgPrime}{blue!12}\colorlet{lgMu}{teal!15}\colorlet{lgLi}{yellow!25}
  \colorlet{lgXi}{gray!22}\colorlet{lgRH}{gray!22}\colorlet{lgGU}{blue!15}\colorlet{lgGL}{red!15}
  \colorlet{lgEArith}{blue!70!black}\colorlet{lgEDef}{black!80}\colorlet{lgEElem}{violet!80!black}
  \colorlet{lgEZero}{gray!45}\colorlet{lgEAnal}{green!50!black}\colorlet{lgEPNT}{brown!70!black}\colorlet{lgEExt}{red!70!black}
  \colorlet{lgTprime}{black}\colorlet{lgTmu}{black}\colorlet{lgTli}{black}
  \colorlet{lgTxi}{gray!50}\colorlet{lgTrh}{gray!50}\colorlet{lgTgamma}{black}
  \colorlet{lgTarith}{black}\colorlet{lgTdef}{black}\colorlet{lgTelem}{black}
  \colorlet{lgTzero}{gray!50}\colorlet{lgTanal}{black}\colorlet{lgTpnt}{black}\colorlet{lgText}{black}
\else
  \colorlet{cPrime}{blue!12}\colorlet{dPrime}{black}\colorlet{tPrime}{black}
  \colorlet{cMu}{teal!15}\colorlet{dMu}{teal!60!black}\colorlet{tMu}{black}
  \colorlet{cLi}{yellow!25}\colorlet{dLi}{orange!70!black}\colorlet{tLi}{black}
  \colorlet{cGammaU}{blue!15}\colorlet{cGammaL}{red!15}\colorlet{dGamma}{black}\colorlet{tGamma}{black}
  \colorlet{cZetaU}{blue!15}\colorlet{cZetaL}{red!15}\colorlet{tZeta}{black}
  \colorlet{cXi}{cyan!15}\colorlet{dXi}{cyan!60!black}\colorlet{tXi}{black}
  \colorlet{cNT}{cyan!15}\colorlet{dNT}{cyan!60!black}\colorlet{tNT}{black}
  \colorlet{cProd}{cyan!15}\colorlet{dProd}{cyan!60!black}\colorlet{tProd}{black}
  \colorlet{cRH}{orange!15}\colorlet{dRH}{orange!70!black}\colorlet{tRH}{black}
  \colorlet{eArith}{blue!70!black}\colorlet{eElem}{violet!80!black}\colorlet{eDef}{black!80}
  \colorlet{ePNT}{brown!70!black}\colorlet{eAnal}{green!50!black}\colorlet{eExpl}{green!50!black}
  \colorlet{eExt}{red!70!black}\colorlet{eZero}{orange!80!black}\colorlet{eLate}{black!80}
  \colorlet{lgPrime}{blue!12}\colorlet{lgMu}{teal!15}\colorlet{lgLi}{yellow!25}
  \colorlet{lgXi}{cyan!15}\colorlet{lgRH}{orange!15}\colorlet{lgGU}{blue!15}\colorlet{lgGL}{red!15}
  \colorlet{lgEArith}{blue!70!black}\colorlet{lgEDef}{black!80}\colorlet{lgEElem}{violet!80!black}
  \colorlet{lgEZero}{orange!80!black}\colorlet{lgEAnal}{green!50!black}\colorlet{lgEPNT}{brown!70!black}\colorlet{lgEExt}{red!70!black}
  \colorlet{lgTprime}{black}\colorlet{lgTmu}{black}\colorlet{lgTli}{black}
  \colorlet{lgTxi}{black}\colorlet{lgTrh}{black}\colorlet{lgTgamma}{black}
  \colorlet{lgTarith}{black}\colorlet{lgTdef}{black}\colorlet{lgTelem}{black}
  \colorlet{lgTzero}{black}\colorlet{lgTanal}{black}\colorlet{lgTpnt}{black}\colorlet{lgText}{black}
\fi\fi
\makeatother

\begin{tikzpicture}[scale=1.0, transform shape,
  box/.style    = {draw, rounded corners=5pt, minimum width=2.8cm,
                   minimum height=0.9cm, align=center, font=\small, thick},
  zbox/.style   = {draw, rounded corners=5pt, minimum width=2.8cm,
                   minimum height=0.9cm, align=center, font=\small,
                   thick, fill=gray!18},
  splitbox/.style = {draw, rounded corners=5pt, minimum width=7cm,
                     minimum height=2.2cm, align=center, font=\small, thick,
                     rectangle split, rectangle split parts=2,
                     rectangle split part fill={red!8, white}},
  arr/.style    = {-{Stealth[length=6pt,width=5pt]}, thick},
  barr/.style   = {{Stealth[length=6pt,width=5pt]}-{Stealth[length=6pt,width=5pt]}, thick},
  % 线型：虚线仅用于 N(T)→RH、ξ↔RH；其余为实线
  proven/.style = {solid},
  cond/.style   = {dashed},
  lab/.style    = {font=\tiny, align=center, inner sep=3pt},
  rhbox/.style  = {draw=dRH, rounded corners=5pt, minimum width=8cm,
                   minimum height=1.0cm, align=center, font=\small,
                   fill=cRH, text=tRH, thick}
]

% 布局与 A-2 同步：左移 1.0cm、上移 2.0cm
% 左右略放宽，避免 Γ/ζ 自环旁公式标注被裁切
\path[use as bounding box]
  (-10.2, -26.2) rectangle (10.0, 7.4);

\begin{scope}[shift={(-1.0,2.0)}]

\node[font=\bfseries\Large, align=center] at (0, 4.35)
  {黎曼猜想知识谱系关系图};

% 主图单独上移 1cm；标题与图例位置不动
\begin{scope}[shift={(-2.75,-1.2)}]
%% ==================== 节点分层 ====================

% 第一层：素数计数函数 (y=0)
\node[box, fill=cPrime, draw=dPrime]    (J)    at (-4.2,  0.0) {$J(x)$};
\node[box, fill=cPrime, draw=dPrime]    (pi)   at ( 0.0,  0.0) {$\pi(x)$};
\node[box, fill=cPrime, draw=dPrime]    (th)   at ( 5.0,  0.0) {$\theta(x)$};
\node[box, fill=cPrime, draw=dPrime]    (ps)   at ( 9.7,  0.0) {$\psi(x)$};

% 第二层：积性函数 (y=-3.0)
% 节点定义修改
\node[box, fill=cMu, draw=dMu] (mu)  at (0.3, -3.0) {$\mu(n)$};
\node[box, fill=cMu, draw=dMu] (Lam) at (5.0, -3.0) {$\Lambda(n)$};

%Weierstrass乘积
\node[box, fill=cProd, draw=dProd, text=tProd]  (prod)   at ( 9.5,  -17.0) {\tiny\text{Hadamard 乘积}};
% M(x) 节点已移除：初等等价信息直接标注在 RH 框内

% 第三层：Li 与 Gamma (y=-6.0 至 -7.0)
\node[box, fill=cLi, draw=dLi] (Li) at (-2.0, -6.0) {$\Li(x)$};

% Gamma(s) 上下分色：上部标收敛域，下部标延拓
\node[splitbox, rectangle split draw splits=false, minimum width=5cm,
      minimum height=4.7cm,
      rectangle split part fill={cGammaU, cGammaL}] (gamma) at (7.0, -6.0)
    {
     \centering \color{tGamma}$\operatorname{Re}(s)>0$
     \nodepart{two}
     \raisebox{6pt}{\color{tGamma}$\mathbb{C}\setminus\{0,-1,-2,\dots\}$}
    };
% Gamma(s) 标注（固定在节点中线左侧）
\node[font=\small, left=-30pt, text=tGamma] at (gamma.west) {$\Gamma(s)$};

% 【修订3】ζ(s) 节点：下半部分补全函数方程左侧 ζ(s)=
\node[splitbox, rectangle split draw splits=false, minimum height=4cm,
      rectangle split part fill={cZetaU, cZetaL}] (zeta) at (0.3, -10.5)
    {\centering\small
     \color{tZeta}$\displaystyle
       \sum_{n=1}^{\infty}\frac{1}{n^s}
       =\prod_{p}\dfrac{1}{1-p^{-s}}
       \quad(\operatorname{Re}(s)>1)$%
     \vphantom{$\displaystyle\sum_{n=1}^{\infty}$}%
     \nodepart{two}
     % 修订：补全 ζ(s)= 使函数方程完整
     \textcolor{tZeta}{$\displaystyle\zeta(s)=
       2^s\pi^{s-1}
       \sin\!\left(\frac{\pi s}{2}\right)\Gamma(1-s)\,\zeta(1-s)$%
     \vphantom{$\displaystyle\sum_{n=1}^{\infty}\mathbb{C}\setminus\{1\}$}}
    };
% ζ(s) 标注


% ξ(s) 节点：下移至与 RH 同一水平，右侧对称于 M(x)
\node[box, fill=cXi, draw=dXi, text=tXi, minimum height=1.2cm]
  (xi) at (4.85, -17.0)
  {$\xi(s)$\\ \tiny 非平凡零点$\rho$与$\zeta(s)$相同\\
  \tiny $\xi(s)=\xi(1-s)$(函数方程推论)
  };


% 【修订8】N(T) 节点：置于 ζ 与 RH 正中间，构成竖直逻辑链
\node[box, fill=cNT, draw=dNT, text=tNT, minimum height=1.3cm]
  (NT) at (2.3, -14.0)
  {$N(T)$\\ \tiny $\zeta(s)$ 零点计数函数\\[1pt]\tiny（幅角原理）};

% ==================== 黎曼猜想节点 ====================
% RH 下移至 y=-20.5，为 N(T) 留出充分空间
\node[rhbox, minimum width=4.6cm, minimum height=2.2cm] (rh) at (-1.9, -17.0) {
    {\tiny\color{tRH} 黎曼猜想 (Riemann Hypothesis, 1859)}\\[3pt]
    {\tiny\color{tRH}$\zeta(s)$ 的所有非平凡零点 $\rho$ 满足
    $\operatorname{Re}(\rho)=\dfrac{1}{2}$}\\[4pt]
{\tiny\color{tRH} 初等等价（梅滕斯函数）：%
  $\displaystyle M(x)=\sum_{n\leq x}\mu(n)$}\\
{\tiny\color{tRH}
  $\mathrm{RH} \iff M(x) = O(x^{1/2+\varepsilon}),\ \forall \varepsilon>0$}
};


% ==================== 连线 ====================

% ----- 第一层内部 -----
\draw[barr, eArith] (J.east) -- (pi.west)
  node[lab, midway, above=14pt, xshift=-2pt, eArith]
    {$\displaystyle\pi(x)=\sum_{n=1}^{\infty}\frac{\mu(n)}{n}J(x^{1/n})$}
  node[lab, midway, below=14pt, xshift=-2pt, eArith]
    {$\displaystyle J(x)=\sum_{n=1}^{\infty}\frac{1}{n}\pi(x^{1/n})$};

\draw[barr, eElem] (pi.east) -- (th.west)
  node[lab, midway, above=14pt, eElem]
    {$\displaystyle\theta(x)=\pi(x)\log x-\int_2^x\frac{\pi(t)}{t}\,\mathrm{d}t$}
  node[lab, midway, above={14pt}, yshift={-50pt}, eElem]
    {\tiny $\displaystyle\pi(x) = \dfrac{\theta(x)}{\log x}+ \int_{2}^{x} \frac{\theta(t)}{t\,(\log t)^2} \,\mathrm{d}t$};

\draw[barr, eElem] (th.east) -- (ps.west)
  node[lab, midway, above=11pt, xshift={5pt}, eElem]
    {$\displaystyle\psi(x)=\sum_{k=1}^{\infty}\theta(x^{1/k})$}
  node[lab, midway, above=-31pt, xshift={5pt}, yshift={-8pt}, eElem]
    {\tiny $\displaystyle\theta(x)=\sum_{k=1}^{\infty}\mu(k)\psi(x^{1/k})$};


% ----- 第二层内部 -----
\draw[barr, eArith] (mu.east) -- (Lam.west)
  node[lab, midway, yshift=20, xshift=0.3cm, eArith]
    {$\displaystyle\Lambda(n)=-\sum_{d\mid n}\mu(d)\log d$
     \quad{\rm (M\"{o}bius 反演)}}
  node[lab, midway, yshift={-12pt-0.5cm}, xshift=0, eArith]
    {$\displaystyle\log n=\sum_{d\mid n}\Lambda(d)$};

% μ(n) → M(x) 的定义连线已省去（见节点内文字），避免与已有连线交叉

% Lambda -> psi
\draw[arr, eDef] (Lam.east) to[out=0, in=-90, looseness=1.0]
  ($(ps.south west)+(30pt,0)$)
  node[lab, pos=0.5, yshift=-90, xshift=250, eDef]
    {$\displaystyle\psi(x)=\sum_{n\le x}\Lambda(n)$};

% ----- 第三层：Li 和 Gamma -----

% 【修订5】Li(x) → π(x)：改为渐近等价语义（非因果推导）
\draw[arr, ePNT]
  ($(Li.north)!.75!(Li.north east)$)
  to[out=90, in=-135, looseness=1.8]
  ($(pi.south west)+(35pt,0)$)
  node[lab, pos=0.4, yshift=-45, xshift=45, ePNT, align=left]
    {\textbf{素数定理（PNT）}：$\pi(x)\sim\Li(x)$\\
     (若RH成立，$\pi(x)=\Li(x)+O(\sqrt{x}\log x)$)};

% Li -> J（显式公式主项；属 eExpl：第2/9章灰白，第15章全文总图着色）
\draw[arr, eExpl]
  (Li.north) to[out=135, in=-90, looseness=1.5]
  ($(J.south east)+(-22pt,0)$)
  node[lab, pos=0.5, yshift=-115, xshift=-65, eExpl, align=center]
    {$\Li(x)$ 是 $J(x)$ 显式\\[2pt]公式的主项};

% Gamma（Re>0）→ ζ：Euler 积分表示（该积分公式本身要求 Re(s)>1）
\draw[arr, eAnal]
  ([yshift=6pt] gamma.west)
  to[out=180, in=90, looseness=0.8]
  ($(zeta.north)+(30pt,0)$)
  node[lab, pos=0.5, yshift=-235, xshift=90,eAnal, align=center]
    {$\displaystyle\zeta(s)\Gamma(s)
      =\int_0^{\infty}\frac{x^{s-1}}{e^x-1}\,\mathrm{d}x$\\
     $(\operatorname{Re}(s)>1)$};

% μ(n) → ζ(s)
\draw[arr, eAnal]
  (mu.south) to[out=-90, in=90, looseness=1.2]
  ($(zeta.north)+({-20pt+0.3cm},0)$)
  node[lab, pos=0.5, yshift={-120pt-0.5cm-0.4cm}, xshift={45pt-5pt}, eAnal]
    {$\displaystyle\frac{1}{\zeta(s)}=\sum_{n=1}^{\infty}\frac{\mu(n)}{n^s}$\\
     $(\operatorname{Re}(s)>1)$};

% ----- J(x) 与 ψ(x) 的 Stieltjes 积分互逆关系 -----
% 【修订4】加注 dψ(t)/ln t 在 t=1 附近需 t>1 的说明
\draw[barr, eArith, line width=0.8pt, rounded corners=8pt]
  ($(J.north)+(0,0.0cm)$)
  -- ++(0,2.0cm)
  -- ++(13.5cm,0)
  -- ++(0,-2.0cm)
  node[lab, pos=0.35, above=6pt, yshift={10pt+0.5cm}, xshift=-200, eArith, align=center]
    {$\displaystyle\psi(x)=\int_{2^-}^x\log t\,\mathrm{d}J(t)$\quad(Stieltjes 积分，$J\to\psi$)}
  node[lab, pos=0.65, below=6pt, yshift={32pt+0.2cm}, xshift=-205, eArith, align=center]
    {$\displaystyle J(x)=\int_{2^-}^x\frac{\mathrm{d}\psi(t)}{\log t}$
     \quad(Stieltjes 反演，$\psi\to J$；积分自 $t>1$)};

% ----- 第四层：ζ 和 ξ 与第二层的连线 -----

% ζ → ψ：显式公式（Perron 积分 + 留数，单向；x≠p^k 时公式精确成立）
\draw[arr, eExpl, rounded corners=30pt]
  ($(zeta.east)+(0, -10pt)$) -| ($(ps.south)+(30pt,0)$)
  node[lab, pos=0.85, right=10pt, yshift=-170, xshift={-170pt-0.2cm},eExpl, align=left]
{$\psi(x)$ 显式公式（Perron 积分，$\zeta\to\psi$）\\
     $\displaystyle \psi(x) = \frac{1}{2\pi i}\int_{c-i\infty}^{c+i\infty}\!\left(-\frac{\zeta'(s)}{\zeta(s)}\right)\!\frac{x^s}{s}\,\mathrm{d}s$\\
     $\displaystyle = x-\sum_{\rho}\frac{x^{\rho}}{\rho}-\log(2\pi)-\tfrac{1}{2}\log(1-x^{-2})$\\
     \tiny $x>1,\;x\neq p^k$；$\rho$ 按 $|\mathrm{Im}(\rho)|$ 升序配对，条件收敛};

% ψ → ζ：Mellin 变换（单向）
\draw[arr, eExpl, rounded corners=30pt]
  ($(ps.south)+(40pt,0)$) |- ($(zeta.east)+(0, -20pt)$)
  node[lab, pos=0.8, yshift=-27, xshift=30,eExpl, align=center]
    {Mellin 变换（$\psi(x) \longleftrightarrow -\dfrac{\zeta'(s)}{\zeta(s)}$，即负对数导数）\\
     $\displaystyle -\frac{\zeta'(s)}{\zeta(s)} = s\int_{1}^{\infty}\frac{\psi(x)}{x^{s+1}}\,\mathrm{d}x$};
       
       
       
       
       
       
       
% Lambda -> zeta（对数微分 Dirichlet 级数）
\draw[arr, eAnal]
  ($(Lam.south west)+(6pt,0)$) to[out=-90, in=90, looseness=0.4]
  ($(zeta.north west)!0.50!(zeta.north east)$)
  node[lab, pos=0.5, yshift={-120pt-0.5cm-0.4cm+1cm}, xshift={145pt-5pt}, eAnal]
    {$\displaystyle-\frac{\zeta'(s)}{\zeta(s)}
      =\sum_{n=1}^{\infty}\frac{\Lambda(n)}{n^s}$};

% J(x) 显式公式（ζ → J）
\draw[arr, eExpl, rounded corners=50pt]
  ($(zeta.south west)+(0,6pt)$) -- ++(-60pt,0) -- ($(J.south west)+(5pt,0)$)
  node[lab, pos=0.5, yshift=-203, xshift={75pt-0.2cm}, eExpl, align=center]
{\textbf{Perron 公式} $\Rightarrow$ $J(x)$ 显式展开\\
     $\displaystyle
     \begin{aligned}
         &J(x)=\frac{1}{2\pi i}\int_{c-i\infty}^{c+i\infty}
           \frac{\log\zeta(s)}{s}\,x^s\,\mathrm{d}s\\
         &=\Li(x)-\sum_{\rho}\Li(x^{\rho})
           -\log 2+\int_x^{\infty}\frac{\mathrm{d}t}{t(t^2-1)\log t}
     \end{aligned}$\\
     \tiny $x>1,\;x\neq p^k$；$\rho$ 按 $|\mathrm{Im}(\rho)|$ 升序配对，条件收敛};

% ζ → J：Mellin 变换（ln ζ）
\draw[arr, eAnal, rounded corners=5pt]
  (zeta.north west) to[out=110, in=-90, looseness=1.5] (J.south)
  node[lab, pos=0.4, yshift=-205, xshift=-55, eAnal, align=center]
    {$\displaystyle\log\zeta(s)=s\int_1^{\infty}J(x)\,x^{-s-1}\,\mathrm{d}x$\\
     $(\operatorname{Re}(s)>1)$};

% ζ → ξ（定义）：ξ 已下移，连线沿右侧空白区域向右下延伸
\draw[arr, eLate, rounded corners=15pt]
  ($(zeta.south east)!0.15!(zeta.north east)$)
  -- ++(1.0, 0)
  -- ++(0, -3.8)
  -- ($(xi.north)+(0pt,0)$)
  node[lab, pos=0.25, right=4pt,yshift=30, xshift=0, eLate, align=center]
    {$\xi(s)=\dfrac{1}{2}s(s-1)\pi^{-s/2}\Gamma(s/2)\zeta(s)$  
};






% xi -> prod (Weierstrass product connection)
\draw[arr, eLate]
  %(xi.south) -- (prod.south)
  (xi.east) --(prod.west)
  node[midway, right=4pt, yshift=28, xshift=-45, eLate]
    {\tiny$\displaystyle\xi(s)=e^{A+Bs}\prod_{\rho}\left(1-\frac{s}{\rho}\right)e^{s/\rho}$};

% prod -> ps (zero sum in explicit formula)
%\draw[arr, violet!70!black, rounded corners=15pt]
  %(prod.east) -- ++(4.4, 0) |- (ps.south)
  %node[lab, pos=0.6, right=2pt, violet!70!black, align=left]
   % {对数微分与 Perron 积分\\
    % 给出显式公式中\\
    % 零点和 $\sum_{\rho}\frac{x_{\rho}}{\rho}$ 项};
%
% ζ 自身左侧弧线（解析延拓）：与表 tab:spectrum-stage8-guide 一致
\draw[arr, eExt, rounded corners=5pt]
  ($(zeta.north west)!0.15!(zeta.south west)$)
  to[out=180, in=180, looseness=1.6]
  ($(zeta.north west)!0.85!(zeta.south west)$);
% 标签用独立节点锚定在 ζ 西侧弧线旁（避免 path 末 node 漂移）
% 单行；中心在 zeta.west 右 2.5cm（相对初版再累计右移）
\node[lab, eExt, align=center, anchor=center]
  at ([xshift=2.5cm]zeta.west)
  {$\zeta$ 亚纯延拓至 $\mathbb{C}\setminus\{1\}$（Hankel / $\vartheta$ 两法）};

% ζ → N(T)：竖直向下（弧线）
\draw[arr, eZero]
  ($(zeta.south east)!0.10!(zeta.north east)$)
  to[out=0, in=90] ++(0.8, -0.8)
  to[out=-90, in=0] (NT.east)
  node[lab, pos=0.6, right=6pt, yshift=-357, xshift=5, eZero, align=center]
    {$\displaystyle N(T)=\frac{T}{2\pi}\log\frac{T}{2\pi e}+O(\log T)$\\
     \tiny 零点计数公式（von Mangoldt）};

% 【已删去】ζ → RH 蓝色直连线（设计意图改为ζ→N(T)→RH线性链）
%\draw[arr, eArith, rounded corners=10pt]
%  (zeta.south)
%  -- ++(0, -3.35)
%  -- ++(-4.1, 0)
%  -- ($(rh.north)!0.80!(rh.north west)$)
%  node[lab, pos=0.35, left=6pt, yshift=20pt,xshift=120,eArith, align=center]
%    {
%     $\zeta(s)$ 的非平凡零点全部位于直线 $\operatorname{Re}(s)=\dfrac12$ 上};
% RH -- P(s) 画线 + 文字
%\draw[arr, cyan!60!black]
 % ($(rh.east)!0.60!(rh.south east)$) -- (prod.west)    % 起点 -- 终点（必须有！）
  %node[midway, right=6pt, font=\tiny, cyan!60!black, align=center]
  %  {$\zeta(s) = \frac{e^{s(\log(2\pi) - 1)}}{2(s-1)} \prod_{n=1}^{\infty} \left(1 + \frac{s}{2n}\right)e^{-s/2n} P(s)$};

% ξ(s) ↔ RH、N(T) → RH：虚线标注（与 RH 相关的谱系连线）

% N(T) → RH：完成 ζ→N(T)→RH 线性逻辑链（虚线）
\draw[arr, eZero, cond, rounded corners=8pt]
  (NT.south west)
  to[out=-150, in=60, looseness=0.9]
  ($(rh.north)+(20pt, 0)$)
node[lab, pos=0.55, left=4pt, yshift=-430, xshift=-25, eZero, align=center]
    {$\mathrm{RH} \iff N_0(T)=N(T),\ \forall\,T>0$\\
   \tiny（$N_0(T)$：临界线 $\operatorname{Re}(s)=\dfrac{1}{2}$ 上虚部 $\leq T$ 的零点数）};

\draw[barr, eZero, cond]
  (xi.west) -- ($(rh.east)!0.0!(rh.south east)$)
  node[lab, midway, above=20pt,yshift=3pt,xshift=5, eZero, align=center]
    {$\mathrm{RH}\iff\xi\!\left(\dfrac{1}{2}+it\right)$\\
     的所有零点 $t$ 均为实数};

% M(x)→RH 箭头已删去：等价信息直接写在 RH 框内

% Gamma 自身解析延拓弧线：与表 tab:spectrum-stage8-guide 一致
\draw[arr, eExt, rounded corners=1.5pt]
  ($(gamma.north east)+(0,-5pt)$)
  to[out=0, in=0, looseness=1.6]
  ($(gamma.south east)+(0,5pt)$);
% 标签用独立节点锚定在 Γ 东侧弧线旁（仅递推式；延拓域已在框下半）
% 相对初版 (gamma.east+14pt)：净向左 3.7cm，上移 2pt
\node[lab, eExt, align=center, anchor=west]
  at ($(gamma.east)+({14pt-3.7cm},2pt)$)
  {递推延拓 $\Gamma(s)=\Gamma(s+1)/s$};

% Gamma（延拓后）→ ζ：Hankel 围道
\draw[arr, eExt]
  ($(gamma.south west)+(40pt,0)$)
  to[out=180, in=0, looseness=1.0]
  ($(zeta.south east)!0.45!(zeta.north east)$)
  node[lab, pos=0.5, yshift=-210, xshift=215, eExt, align=center]
    {$\displaystyle\zeta(s) = -\frac{\Gamma(1-s)}{2\pi i}\oint_{\mathrm{H}}\frac{(-z)^{s-1}}{e^z-1}\,\mathrm{d}z$\\
     解析延拓到 $\mathbb{C}\setminus\{1\}$，Hankel 围道消去分支切割};

\end{scope}% 主图

% ==================== 图例（取自 A-2 详细文案 + 分阶段灰白色） ====================
\node[draw=black!40, fill=gray!8, rounded corners=8pt, thick, align=center, inner sep=10pt, font=\tiny]
  (legend) at (0, -23.5) {
  \begin{tabular}{@{}ll@{\hspace{1.0cm}}ll@{}}
    \multicolumn{2}{l}{\scriptsize\textbf{节点颜色分类图例}} &
    \multicolumn{2}{l}{\scriptsize\textbf{知识谱系连线语义图例}} \\[6pt]
    \tikz[baseline=-0.5ex]\fill[lgPrime, draw=black!70, thick, rounded corners=1.5pt] (0,-0.12) rectangle (0.55,0.12); &
    \textcolor{lgTprime}{\textbf{素数计数函数}（实变，如 $\pi(x)$）} &
    \textcolor{lgEArith}{\rule[0.5ex]{16pt}{1.4pt}} &
    \textcolor{lgTarith}{\textbf{算术互逆}（M\"{o}bius / Stieltjes 反演）} \\[3pt]
    \tikz[baseline=-0.5ex]\fill[lgMu, draw=black!70, thick, rounded corners=1.5pt] (0,-0.12) rectangle (0.55,0.12); &
    \textcolor{lgTmu}{\textbf{算术积性函数}（离散，如 $\mu(n)$）} &
    \textcolor{lgEDef}{\rule[0.5ex]{16pt}{1.4pt}} &
    \textcolor{lgTdef}{\textbf{对象定义}（构造定义与主项提取）} \\[3pt]
    \tikz[baseline=-0.5ex]\fill[lgLi, draw=black!70, thick, rounded corners=1.5pt] (0,-0.12) rectangle (0.55,0.12); &
    \textcolor{lgTli}{\textbf{渐近逼近主项}（如 $\Li(x)$）} &
    \textcolor{lgEElem}{\rule[0.5ex]{16pt}{1.4pt}} &
    \textcolor{lgTelem}{\textbf{初等计数变换}（数论函数基本恒等式）} \\[3pt]
    \tikz[baseline=-0.5ex]\fill[lgXi, draw=black!55, thick, rounded corners=1.5pt] (0,-0.12) rectangle (0.55,0.12); &
    \textcolor{lgTxi}{\textbf{辅助复变函数}（如 $\xi(s)$）} &
    \textcolor{lgEZero}{\rule[0.5ex]{16pt}{1.4pt}} &
    \textcolor{lgTzero}{\textbf{零点结构与猜想}（分布特征及等价命题）} \\[3pt]
    \tikz[baseline=-0.5ex]\fill[lgRH, draw=black!55, thick, rounded corners=1.5pt] (0,-0.12) rectangle (0.55,0.12); &
    \textcolor{lgTrh}{\textbf{核心猜想与目标}（如 RH）} &
    \textcolor{lgEAnal}{\rule[0.5ex]{16pt}{1.4pt}} &
    \textcolor{lgTanal}{\textbf{分析桥梁}（积分变换与 Dirichlet 级数）} \\[3pt]
    \tikz[baseline=-0.5ex]{\clip[rounded corners=1.5pt] (0,-0.12) rectangle (0.55,0.12); \fill[lgGL] (0,-0.12) rectangle (0.55,0); \fill[lgGU] (0,0) rectangle (0.55,0.12); \draw[black!55, thick, rounded corners=1.5pt] (0,-0.12) rectangle (0.55,0.12); \draw[black!55, thick] (0,0) -- (0.55,0);} &
    \textcolor{lgTgamma}{\textbf{解析延拓复变}（初始域/延拓域）} &
    \textcolor{lgEPNT}{\rule[0.5ex]{16pt}{1.4pt}} &
    \textcolor{lgTpnt}{\textbf{渐近定理}（PNT 等确立的渐近推论）} \\[3pt]
    & & \textcolor{lgEExt}{\rule[0.5ex]{16pt}{1.4pt}} &
    \textcolor{lgText}{\textbf{解析延拓}（复围道积分与全纯延拓）} \\[3pt]
    & & \tikz[baseline=-0.5ex]{\draw[lgEZero, line width=1.2pt] (0.00,0) -- (0.12,0); \draw[lgEZero, line width=1.2pt] (0.17,0) -- (0.29,0); \draw[lgEZero, line width=1.2pt] (0.34,0) -- (0.46,0);} &
    \textbf{未证}
  \end{tabular}
};

\end{scope}% 位移：左 1.0cm、上 2.0cm

\end{tikzpicture}
%
}
\end{figure}
\clearpage
\noindent
\captionof{figure}{黎曼猜想知识谱系关系图（第二次着色：至函数方程）。
与图~\ref{fig:spectrum-stage2}（第一次）、图~\ref{fig:function-relations}（第三次）同一底图；
彩色为截至本章已建立的部分，灰白为后文。
相对第一次进度新增的彩色节点/连线与正文对照见表~\ref{tab:spectrum-stage8-guide}。}
\label{fig:spectrum-stage8}

\paragraph{图~\ref{fig:spectrum-stage8} 新增彩色元素与正文对照}
下列各项将图上\textbf{相对第一次进度图新着色}的节点与连线引向相应定义与证明
（含本章及前序已建立的 $\operatorname{Re}(s)>1$ 分析桥梁）；
初等层（$\pi,J,\theta,\psi,\mu,\Lambda,\Li$ 及其算术互逆）见
表~\ref{tab:spectrum-stage2-guide}，此处不重复。
\begin{table}[htbp]
\centering
\caption{图~\ref{fig:spectrum-stage8} 新增彩色图元与正文对照}
\label{tab:spectrum-stage8-guide}
\small
\renewcommand{\arraystretch}{1.35}
\begin{tabular}{@{}p{4.9cm}p{8.3cm}@{}}
\toprule
\textbf{图上元素} & \textbf{正文（定义 / 公式 / 证明）} \\
\midrule
节点 $\zeta(s)$（上半：$\operatorname{Re}s>1$） &
  第~\ref{ch:riemann_zeta_intro}~章：Dirichlet 级数与 Euler 乘积定义；
  第~\ref{sec:holo}~节：$\operatorname{Re}s>1$ 上的全纯性 \\
节点 $\zeta(s)$（下半：解析延拓） &
  第~\ref{sec:elementary_extension}~节：$\eta$ 交错级数延拓至 $\operatorname{Re}s>0$
  （式~\eqref{eq:zeta_eta_continuation}）；
  第~\ref{sec:riemann_first_method}~、\ref{sec:theta_function_method}~节：
  全平面亚纯延拓（仅 $s=1$ 为极点） \\
节点 $\Gamma(s)$（上半：$\operatorname{Re}s>0$） &
  第~\ref{ch:gamma_function}~章第~\ref{sec:gamma_def}~节：Euler 积分定义 \\
节点 $\Gamma(s)$（下半：解析延拓） &
  第~\ref{sec:gamma_continuation}~节：递推延拓至 $\mathbb{C}\setminus\{0,-1,-2,\ldots\}$；
  第~\ref{sec:reflection_formula}~节：余元公式 \\
$\zeta$ 节点左侧自环（解析延拓） &
  第~\ref{sec:riemann_first_method}~节式~\eqref{eq:zeta_contour_gamma}、
  第~\ref{sec:theta_function_method}~节式~\eqref{eq:analytic_continuation}：
  Hankel / $\vartheta$ 两法给出的全局延拓 \\
$\Gamma$ 节点右侧自环（解析延拓） &
  第~\ref{sec:gamma_continuation}~节定理~\ref{thm:gamma-meromorphic} \\
连线 $\Gamma\to\zeta$（Euler 积分，
  $\zeta(s)\Gamma(s)=\int x^{s-1}/(e^x-1)\,\mathrm{d}x$） &
  第~\ref{ch:riemann_zeta_intro}~章第~\ref{sec:integral-representation-eq}~节\eqref{eq:zeta-integral}；
  第~\ref{ch:gamma_function}~章缩放式\eqref{eq:gamma-scaled}；
  本章第~\ref{sec:riemann_first_method}~节由此出发 \\
连线 $\Gamma\to\zeta$（Hankel 围道，
  $\zeta=-\Gamma(1-s)/(2\pi i)\oint_{\mathrm{H}}\cdots$） &
  第~\ref{sec:riemann_first_method}~节：
  式~\eqref{eq:zeta_contour_gamma} 与路径变形；
  同节导出函数方程~\eqref{eq:asymmetric_fe}、\eqref{eq:asymmetric_fe_alt} \\
连线 $\mu\to\zeta$
  （$1/\zeta(s)=\sum\mu(n)n^{-s}$，$\operatorname{Re}s>1$） &
  第~\ref{ch:riemann_zeta_intro}~章
  第~\ref{subsec:arith_dirichlet_gen}~节 \\
连线 $\Lambda\to\zeta$
  （$-\zeta'/\zeta=\sum\Lambda(n)n^{-s}$，$\operatorname{Re}s>1$） &
  第~\ref{ch:riemann_zeta_intro}~章
  第~\ref{subsec:chebyshev_zeta_relation}~节 \\
连线 $\zeta\to J$（Mellin，
  $\log\zeta(s)=s\int_1^\infty J(x)\,x^{-s-1}\,\mathrm{d}x$） &
  第~\ref{ch:integral_transforms}~章：
  Dirichlet 级数与部分和的 Mellin 联系
  （$F(s)=s\int A(x)x^{-s-1}\,\mathrm{d}x$）；
  系数 $A=J$ 时 $F=\log\zeta$（$\operatorname{Re}s>1$） \\
第二法 $\vartheta$ 与对称积分表示
  （图上函数方程所依赖的另一路径） &
  $\vartheta$ 的定义与模反演（Gauss 型 FT + Poisson）
  见第~\ref{ch:integral_transforms}~章第~\ref{subsec:gauss_ft}~节及随后的模反演小节；
  应用于 $\zeta$ 见第~\ref{sec:theta_function_method}~节式~\eqref{eq:analytic_continuation}；
  与第一法对照见表~\ref{tab:riemann_methods_comparison}
  （第~\ref{sec:riemann_two_methods_comparison}~节） \\
平凡零点 $\zeta(-2k)=0$
  （函数方程推论，左半平面几何） &
  第~\ref{sec:critical_strip_from_fe}~、\ref{sec:trivial_zeros_geometry}~节
  与图~\ref{fig:trivial_zeros}；由~\eqref{eq:asymmetric_fe_alt} 推出 \\
\bottomrule
\end{tabular}
\end{table}
\begin{center}
\fbox{\parbox{0.92\textwidth}{
\centering
\vspace{0.2em}
\textbf{第九章：重要定理与核心公式汇总}\\[6pt]
\renewcommand{\arraystretch}{1.6}
\hbadness=10000
\begin{tabular}{lp{9.5cm}}
\hline
\textbf{名称} & \textbf{数学内容 / 公式表达} \\
\hline
交错级数关系式 & $\zeta(s) = \dfrac{\eta(s)}{1-2^{1-s}}, \quad \Re(s)>0, \ s \in \mathbb{C} \setminus \{1\}$ \\
第一法围道积分表示 & $\displaystyle\zeta(s) = -\dfrac{\Gamma(1-s)}{2\pi i}\displaystyle\int_C \dfrac{(-z)^{s-1}}{e^z-1}dz, \quad s \in \mathbb{C} \setminus \{1\}$ \\
第二法积分表示 & $\pi^{-s/2} \Gamma(s/2) \zeta(s) = \dfrac{1}{s(s-1)} + \displaystyle\int_1^\infty \left(x^{s/2-1} + x^{-(s+1)/2}\right)\omega(x)dx, \quad s \in \mathbb{C} \setminus \{0, 1\}$ \\
非对称函数方程 & $\zeta(s) = 2^s \pi^{s-1}\sin\left(\dfrac{\pi s}{2}\right)\Gamma(1-s)\zeta(1-s), \quad s \in \mathbb{C} \setminus \{1\}$ \\
平凡零点定理 & 当 $s = -2k$ ($k \in \mathbb{Z}^+$) 时，恒有 $\zeta(-2k) = 0$ \\
\hline
\end{tabular}
\vspace{0.2em}
}}
\end{center}
